% 本文件由 LaTeX 排版 Agent 根据有效 translation.json 自动生成。
% author=Gregory of Nyssa; work_id=2915; title=论灵魂与复活

\chapter{论灵魂与复活}

% structure: p0001 source=h1 target=omitted reason=page-title-already-rendered-by-parent

在圣人中伟大的巴西略已离开此世，归向天主；各教会共同为他哀恸。但他的妹妹，那位女教师，仍活在世上；于是我前往她那里，渴望因她失去兄弟而交换同情。我的灵魂因这沉重打击而极度悲伤，我寻找一个能同样感受此事的人，好与他同洒眼泪。但当我们彼此相见时，女教师的容颜唤醒了我所有的痛苦；因为她也正躺卧在病榻上，几乎奄奄一息。她如同一位熟练的驭手，在我无法控制的悲痛中，先让我稍作宣泄；然后她试图用言语来制止我，并用她推理的缰绳纠正我灵魂的混乱。她引用了宗徒的话，关于不应\BibleQuote{为安眠的人忧伤}的责任；因为只有\BibleQuote{没有希望的人}才有这样的情感。我的心中仍在翻腾着痛苦，于是我问道——\par

\Quote{人类如何能实践这一点？对死亡的憎恶是如此本能地深植于一切人心中！那些看着临终之床的人几乎无法忍受那景象；而那些死亡临近的人则尽其所能地退缩。为什么，甚至管制我们的法律也将死刑列在罪行之首，也列在刑罚之首。那么，我们用什么方法能使自己视离开生命为无物，哪怕是陌生人的死亡，更不用说亲人了，当他们停止生存之时？我们看到整个人生历程都指向这一件事，即我们如何能继续活在此生；事实上，正是为此我们发明了房屋来居住；使我们的身体不致因寒冷或炎热而在环境中倒下。农业，又是什么？不过是提供我们的生计？实际上，一切关于我们如何继续活下去的思考都是由对死亡的恐惧所引起的。为什么医学在人们中如此受尊崇？因为人们认为它通过其方法在某种程度上与死亡作战。为什么我们有胸甲、长盾、胫甲、头盔，以及所有防御性盔甲，还有堡垒的围墙，和铁栅门？不正是因为我们害怕死亡吗？死亡既然对我们如此自然地可怕，那么幸存者又如何能轻易地遵守这项命令，对逝去的朋友无动于衷呢？}\par

女教师问道：\Quote{那么，仅仅在死亡本身的必然性中，你感到什么特别的痛苦？那些无思想之人的通常说法不足以构成指控。}\par

我回答她：\Quote{怎么！难道没有悲恸的缘由吗？当我们看见一个刚才还在生活说话的人，突然变得毫无生气、一动不动，每个身体器官的感觉都消失了，没有视觉或听觉在运作，也没有任何感觉所拥有的其他领会能力；如果你用火或铁触碰他，即使你拔剑刺入身体，或将其抛给野兽，或将其埋在土丘之下，那死人面对任何处置都同样无动于衷？那么，既然这种变化以所有这些方式被观察到，而生命的原理——无论它是什么——突然从眼前消失，就像一盏熄灭的灯，刚才还燃烧着的火焰既不留在灯芯上，也不转移到其他地方，而是完全消失，那么，一个没有明确根据的人，怎能不带着情感承受这样的变化？我们\Emph{听见}灵魂的离去，我们\Emph{看见}留下的躯壳；但对于那已分离的部分，我们一无所知，无论是它的本质，还是它逃往何处；因为无论是地、空气、水，还是其他任何元素，都无法表明这个离开身体的力驻留在自身之内；一旦它离开，只留下一具尸体，等待分解。}\par

当我这样滔滔不绝地讲述时，女教师向我示意，说道：\Quote{难道使你惊慌不安的，不是灵魂不会永远持续，而是随着身体的解体而终止的想法吗？}\par

我有点大胆地回答，并没有仔细考虑我的话，因为我的悲痛情感尚未让我恢复判断力。实际上，我说：\Quote{在我看来，神圣的宣告就像一些命令，强迫我们相信灵魂永远持续；然而，我们并非由任何推理而引导至这种信仰。我们内在的心灵似乎奴性地接受这种强加的意见，但并不以自发的冲动而默许。因此，我们为逝者而悲伤更加沉重；我们并不确切知道这个赋予生命的原理是否本身是独立的；它在哪里，或如何存在；事实上，它是否在任何地方以任何方式存在。这种对实际情况的不确定性使两边的意见平衡；许多人采纳一种观点，许多人采纳另一种观点；而且确实有某些在希腊哲学家中享有不小名声的人，他们持有并坚持了我刚才所说的。}\par

她喊道：\Quote{摒弃那种异教的无稽之谈！因为谎言之父在其中捏造虚假的理论，只为危害真理。你只须注意这一点：这种关于灵魂的观点无异于放弃德行，只追求当下的快乐；永生——唯独凭借它，德行才能声称其优势——就必须绝望。}\par

我问道：\Quote{那么，我们如何能对灵魂的延续获得坚定而不可动摇的信仰？我也意识到，人类生命将失去生命所能给予的最美的点缀——我指德行——除非我们内心对此建立起无疑的信心。对那些认为此生是生存的极限、不盼望任何死后之事的人，德行究竟有何立足之地？}\par

女教师回答说：\Quote{那么，我们必须寻找从何处开始我们关于这点的讨论；如果你愿意，由你自己承担对立观点的辩护；因为我看你的心思有点倾向于接受这样的辩护。然后，在陈述了相冲突的信念之后，我们就能寻找真理了。}\par

当她提出这个请求时，我打消了她真以为我要提出反对意见的疑虑，我并非真的想反对，而只是想先传唤那些反对灵魂信仰的观点来审视，以此为信仰找到一个坚实的基础。于是，我开始说道——\par

持相反信仰的人难道不会这样说吗：身体作为组合体，必然要分解成其构成成分？当身体中各元素的联合停止时，每个元素都自然地趋向于与自己同类的元素，带着一种不可抗拒的相近相吸的倾向；我们体内的热将与热结合，土性的部分与固体结合，其他每个元素也将移向自己的同类。那么，灵魂此后会在哪里呢？如果有人肯定灵魂就在这些元素中，他就不得不承认灵魂与这些元素是同一的，因为如此不同性质的两样东西之间不可能发生融合。但这样一来，灵魂必然被视为一个复杂的东西，与如此对立的性质融合在一起。然而复杂并非单一，必须被归类为复合物，而复合物必然是可分解的；分解意味着复合物的毁灭；可毁灭之物不是不朽的，否则连肉身本身，既然可分解为构成元素，也可以被称为不朽了。另一方面，如果灵魂是不同于这些元素的某种东西，那么当灵魂由于异质的本性既不能在这些元素中被发现，而在世界上也没有其他地方可以按其特有性质继续存在时，我们的理性又能指出一个地方让灵魂在哪里呢？如果一样东西无处可寻，显然它就不存在。\par

导师听到我这些话，轻轻叹了口气，然后说道：也许正是这些反对意见，或者诸如此类的说法，是斯多葛派和伊壁鸠鲁派在雅典收集起来反驳宗徒的。我听说伊壁鸠鲁正是将他的理论引向这个方向。在他看来，万物的结构是偶然的、机械的，没有一个 \OriginalTerm{天主眷顾}{Providence} 贯穿其运作；与此一致，他认为人的生命就像一个气泡，只要内在的气息被包裹的物质所维系，就存在片刻，因为我们的身体只是一个膜，包裹着一口气；一旦膨胀消散，被囚禁的本质就熄灭了。对他而言，可见之物就是存在的界限；他使我们的感官成为我们认识事物的唯一手段；他完全闭上了灵魂的眼睛，无法看到理性和非物质世界中的任何事物，就像一个被囚禁在小屋里的人，墙壁和屋顶遮挡了外界视线，无法看到天空的一切奇观。的确，宇宙中凡是被视为感觉对象的一切，都如同土墙，在较狭隘的灵魂与那等待着它们默观的理性世界之间形成了一道屏障；这样的灵魂只能看见土、水、火；这些元素从何而来，在什么之中、被什么包围，这些灵魂因其狭隘无法察觉。而看到一件衣服，任何人都会想到织工；看到船，就想到造船者；看到建筑，就想到建造者的手；这些狭隘的灵魂凝视世界，但他们的眼睛对那通过我们周围万物彰显自身的那一位是盲目的；于是他们提出了那些关于灵魂消亡的机巧而尖锐的学说：身体来自元素，元素来自身体，此外，灵魂不可能自存（如果它不是这些元素之一，或不寄居在某个元素中）；如果这些反对者认为，由于灵魂与元素本性不同，\Emph{死后}便无处存身，那么他们首先必须提出，灵魂在有肉身的生活中也不存在，因为身体本身不过是这些元素的集合；因此他们不能告诉我们灵魂在那里独立地赋予它们的组合物以生命。如果死后灵魂不可能存在，而元素却存在，那么，我说，按照这种教导，我们的生命也证明无非是死亡。但另一方面，如果他们不怀疑现在灵魂在身体中的存在，那么当身体分解为元素时，他们怎能坚持灵魂的消亡呢？然后，其次，他们必须以同样的胆量针对这种本性中的天主。因为他们怎能断言那理性和非物质、不可见的对象能被溶解并扩散到湿与软、热与干之中，从而通过渗透万物，虽与所渗透之物本性不同，却并非不能渗透，以此维系宇宙的存在呢？因此，让他们从他们的体系中除去那支撑世界的天主本身吧。\par

我说道：这正是我们的对手必定会怀疑的一点；即万物都依赖天主并由祂所包围，或者有任何超越有形世界的天主存在。\par

童贞女喊道：\Quote{对于这样的疑虑，沉默更为合适，并且不屑于对如此愚蠢和邪恶的命题作出任何回答；因为有一条神圣的诫命禁止我们按愚昧人的愚昧回答他；正如先知所宣称的，凡说\InnerQuote{没有天主}的人必定是愚昧人。但既然必须说话，我要向你们提出一个论点，这论点不是我的，也不是任何人的（因为若是人的，无论谁说的，价值都很小），而是整个受造界通过它的奇妙，以精巧而富有艺术性的、深入心灵的话语，向眼睛的听众所宣告的论点。受造界直接宣告了造物主；因为诸天，正如先知所说，以它们不可言喻的话语宣扬了天主的荣耀。我们看见这奇妙的天和奇妙的地中的普遍和谐：本质上彼此对立的元素如何以一种不可言喻的合一交织在一起，服务于一个共同的目的，每个元素贡献其特定的力量来维持整体；那些不相混和、相互排斥的元素如何不因它们的特性而彼此分离，也不因这样的对立在混合时被毁灭；那些天生轻飘的元素如何向下运动，例如太阳的热量在光线中下降，而沉重的物体却通过稀薄化为蒸汽而上升，以致水违反其本性上升，通过空气被输送到上层区域；穹苍的火如何渗透大地，甚至其深渊也感到热量；注入土壤的雨水湿气如何，虽然本性为一，却产生无数不同的胚芽，并按适当比例赋予每个受其影响的主题以生气；极地天球如何非常迅速地旋转，其内部的轨道如何反向运动，以及所有行星的食、合和测定的间隔。我们用心灵的锐利眼睛看见这一切，并且无法不通过这样的景象被教导：一种神圣的能力，以技巧和秩序运作，在这个有形世界中彰显自己，渗透每个部分，将这些部分与整体结合，并通过部分完成整体，以单一的、自控自足、从不停止运动、却永不改变其所处位置的能力环绕宇宙。}\par

我问道：\Quote{请问，这种对天主存在的信仰怎么会同时证明人类灵魂的存在呢？因为天主当然与灵魂不是同一的，以致如果相信了其中一个，另一个就必然被相信。}\par

童贞女回答说：\Quote{有智慧的人曾说过，人是一个小世界，在他自己内包含着构成宇宙的所有元素。如果这个观点是真的（而且看来如此），我们也许就不需要别的帮手来确立我们关于灵魂的概念的真理了。我们对它的概念是这样的：它存在，具有一种罕见而独特的本性，独立于具有粗糙质地的身体。我们是通过感官的把握来获得关于这个外部世界的确切知识的，而这些感官活动本身又引导我们去理解超感觉的事实与思想的世界，我们的眼睛因而成为那种在宇宙中可见的全能智慧的诠释者，并在自身中指向那环绕它的存在。同样，当我们审视我们的内心世界时，我们在已知的事物中也找到了不少线索来推测未知的事物；那里的未知事物也正是那作为思想对象而非视觉对象、逃避感官掌握的东西。}\par

我回答说：\Quote{不，从在这和谐的物质自然构造中可观察到的巧妙而艺术的设计，很可能推断出一种超越宇宙的智慧；但是，就灵魂而言，对于那些想要从身体所提供的任何迹象，通过已知来追溯未知的人，有什么知识是可能的呢？}\par

童贞女回答说：\Quote{毫无疑问，对于愿意遵循智慧箴言\InnerQuote{认识你自己}的人，灵魂本身就是一位称职的导师；我的意思是，她是一个非物质的和精神的事物，以与她独特本性相应的方式运作和运动，并通过身体的器官显明这些独特的情感。因为即使是在那些刚因死亡而成为尸体的人身上，这个身体组织也存在着，但它没有运动或行动，因为灵魂的力量不再在其中。只有当感觉在器官中存在时，它才运动，而且不仅如此，精神的力量通过那种感觉，带着它自己的冲动渗透并随意运动所有那些感觉器官。}\par

我问：\Quote{那么，灵魂是什么呢？也许可能有某种方法描绘其本性，以便我们能够以草图的方式对这个主题有所理解。}\par

童贞女回答说：\Quote{不同的作者曾以不同的方式尝试过它的定义，各人根据各自的倾向；但以下是我们关于它的看法。灵魂是一种受造的、有生命的、理智的本质，它从一个有组织的、有感觉的身体传递生命和把握感觉对象的能力给自身，只要一个能够这样做的自然构成保持合一。}\par

她说这话时指着坐在旁边观察她病情的医生，说：\Quote{我们身边就有一个证明我所说之事的证据。}请问，这个人如何通过将手指搭在脉搏上，借着触觉在某种意义上听到了自然大声向他呼喊并告知她特有的痛苦？事实上，他得知体内的疾病是炎症性的，病源出自某个内脏器官，以及发热的程度如何？他又如何通过眼睛的观察被教导其他这类事实，当他查看病人的姿势、观察肌肉的消瘦时？同样，面色苍白而带胆汁质、眼神如疼痛者不由自主地倾向于悲伤，这些也指示了内在的状况；而耳朵则通过呼吸短促和随之而来的呻吟了解病情的性质，从而获得类似的信息。可以说，即使是专家的嗅觉也并非不能察觉失调的种类，而是通过呼吸的特定气味注意到内脏的隐秘痛苦。如果每个感官器官中没有某种智能的力量，这一切怎能发生？没有思想引导它从感觉到理解面前的对象，我们的手本身能教导我们什么？没有心灵，耳朵、眼睛、鼻孔或任何其他器官本身，对解决问题有何帮助？诚然，一位外教学者曾说过一句极为真实的话：\Quote{是心灵在看，是心灵在听。}否则，如果你不承认这是真的，你必须告诉我，为什么当你按照导师教导的方式观察太阳时，你断言太阳的圆面并非如众人所见的那样大，而是比整个地球的尺寸大许多倍？你不是确信如此，因为你通过推理从现象出发，得出了关于某种运动、如此的时间和空间距离、以及日食的原因的认识吗？当你观察月亮的亏缺和盈满时，天体的可见形象教导你其他真理，例如：月亮本身没有光，它在最靠近地球的轨道上运行，并且借太阳的光照亮；就像镜子一样，它接受太阳光，反射的不是自己的光线，而是太阳的光线，这些光线从它光滑闪烁的表面反射回来。看见这些而不仔细考察的人，以为光来自月亮本身。但事实并非如此，可由以下证明：当月亮正对太阳时，它朝向我们的整个圆面都被照亮；但由于它在更小的空间内运行，其公转速度更快，在太阳绕行一周之前，月亮已经完成了十二次以上的公转；因此，它的实体并非始终被光覆盖。因为在她频繁的公转中，正对太阳的位置无法保持；当这个位置导致朝向我们的整个月面被照亮时，一旦她侧移，朝向我们的半球必然部分被阴影遮盖，只有朝向太阳的部分接受它拥抱的光线；事实上，光明持续从不再见到太阳的部分退却到仍见到太阳的部分，直到她完全越过太阳的圆面，在其背面接受太阳光线；这时，她本身完全没有光和辉耀的事实使得朝向我们的侧面不可见，而另一个半球则全在光中；这被称为亏缺的完成。但是，当她再次在自己的公转中经过太阳并横对太阳光线时，刚才黑暗的侧面开始略微发光，因为光线从被照明的部分移动到不久前还不可见的部分。你看，眼睛的确教导了这些；然而，如果没有某种通过眼睛观看并利用感官数据作为向导从显见之物穿透到不可见之物的东西，眼睛本身永远无法提供这种洞见。无需再补充几何学的方法，它通过可见的图形逐步引导我们到达不可见的真理，以及无数其他例子，这些都证明领会是我们本性深处的一种理智实体的工作，它通过我们身体感官的运作而行动。\par

但是如果有人坚持认为，尽管所有元素在其可见形式中共享某种物质性质，但每种特定物质之间存在巨大差异（例如，运动并非在所有元素中相同，有些向上运动，有些向下；形式和质量也是如此），并且说同样有一种力量以纯粹自然的方式被纳入和属于这些元素，从而实现这些理智的洞察和运作（例如，我们常常在机械师的作品中看到类似的效果，他们按照技艺规则组合物质，从而模仿自然，不仅在形态上相似，甚至在运动上也相似，以至于当机械装置在发声部分发出声音时，它模仿人的声音，然而我们无法在任何地方感知到有某种心智力量在制作特定的形态、性格、声音和运动）；假设我说，我们断言这一切在我们自然身体的有机机器中同样产生，没有混合任何特殊的思考实体，而仅仅由于元素本身内在的推动力自行完成这些运作——实际上，无非是一种朝向认识面前对象的冲动运动；那么，难道这不正证明了那种你所说的理智而不可触知的存在——灵魂——绝对不存在吗？\par

她回答说：\Quote{你的例子和据此所作的推论，虽然属于反驳论证，但都可以成为我们陈述的盟友，并且将大大有助于证实它的真理。}\par

为什么，你怎能这样说？\par

因为，你看，如此理解、操控和处置没有灵魂的物质，以至于储存在这类机械中的技艺几乎成为这物质的灵魂，在各种方式上模仿运动、形状和声音等，可被转化为一种证明：在人内存在着某种东西，使他显示一种内在的能力，通过沉思和发明才能，在自己内部想出这样的构思，并在理论上准备这类机械后，通过手工技巧付诸实践，在物质中展现其心智的产物。例如，首先他通过思考看出，要产生任何声音都需要某种风；然后，为了在机械中产生风，他事先通过推理和对元素性质的仔细观察确定，世界上根本没有真空，而只是与较重的相比，较轻的才被认为是真空；因为空气本身，作为一种独立的存在，是拥挤而充实的。说一个罐子是\Quote{空的}，是语言的滥用；因为当它没有任何液体时，即使在这种状态下，对有经验的人来说也仍然是满的。有一个证明：当一个罐子放入水池中时，它不会立即充满，而是先浮在水面上，因为它所含的空气帮助浮起其圆面；直到最后打水者的手将其压到水底，在那里它通过颈部进水；在这个过程中，即使在水进入之前，它也显示出并非空的；因为可以观察到在颈部两种元素之间发生一场争斗：水因其重量被迫进入内部，因此涌入；而被囚禁的空气则因水沿颈部涌入而空间受限，于是反向冲出；因此水被空气的强流阻挡，发出汩汩声和气泡声。人们观察到此，并根据这两种元素的这一特性设计出一种引入空气以驱动机械的方法。他们用某种硬质材料制成一个空腔，阻止其中的空气向任何方向逸出；然后通过其口将水引入这个空腔，根据所需量分配水量；接着允许空气向相反方向出口，使其进入一个准备好的管子，在此过程中被水强力挤压而成为一股气流；这气流在管子的结构上吹奏，产生一个音符。难道这些可见的结果不是清楚地证明，在人内存在着某种心智，一种不同于可见物的东西，它凭借自身不可见的思维本性，首先通过内在的发明预备这样的装置，然后在它们成熟后，将其带到光明中并在顺从的物质中展现出来？因为如果有可能像我们对手的理论那样将这些奇迹归功于元素的实际构造，那么这些机械就会自行建造；青铜不会等待艺术家来做成人的形状，而会凭借内在力量成为那样；空气不需要管子来发出音符，而会凭借自身随机的流动和运动自发发声；水向上喷射也不会像现在这样是人为压力迫使它沿非自然方向运动的结果，而是水会自行进入机械，在该方向找到自然通道。但如果这些结果中没有一个是自发由元素力产生的，而是相反，每种元素都被技巧随意使用；而技巧是心智的一种运动和活动，那么由反对观点所提出的这些后果本身，岂不向我们显示心智是某种不同于被感知的东西吗？\par

那被感知的东西，我回答说，不是同一回事，我承认；但我并未从这样的陈述中发现对我们问题的任何答案；我还不清楚我们应当认为那不被感知的东西是什么；你的论证只向我显示它不是任何物质的东西；我仍不知道它合适的名字。我特别想知道它是什么，而不是它不是什么。\par

我们确实通过同一种方法学习到许多关于许多事物的知识，她回答说，因为就在说一件事物是\Quote{不是这样那样的}时，我们通过暗示解释那事物的本性本身。例如，当我们说\Quote{无欺诈的}时，我们指一个好人；当说\Quote{无丈夫气的}时，我们表达了那人是个懦夫；同样可以暗示许多其他事物，其中我们或通过否定坏来表达好，或反之\Emph{vice versâ}。因此，如果有人对当前的主题这样想，他就不会错失对它的正确概念。问题是——我们应当如何认为心智在其本质本身中？既然发问者因它向我们显示的活动而对所问之事的存在已无疑问，只想知其是什么，那么被告知它不是我们感官所感知的，既不是颜色，也不是形状，也不是硬度，也不是重量，也不是数量，也不是立方尺寸，也不是点，也不是任何在物质中可感知的东西，他就会充分发现它是什么了——假设确实存在一个超出这一切的东西。\par

此时我打断了她的论述说：你若把这些全部排除，我不明白如何能避免连你所寻求的那个东西也一并取消。就目前而言，我无法设想，离开了这些，知觉活动还能依附于什么。因为每当用审视的理智考察世界的内容时，我们只要着手触及所寻求之物，就必然像盲人沿墙壁摸索门扉一样，碰到上述事物中的某一个；我们必定遇到颜色、形状、数量，或你列举的某样东西；而如果说那事物并非这些中的任何一个，我们心智的软弱便会促使我们以为它根本不存在。\par

她愤慨地打断说：这种谬论真可耻！这狭隘鄙陋的世界观竟将我们带到何等结论！若一切感官不可认知之物都要从存在中被抹去，那么同一声明便承认那统摄万物的无所不包的大能不存在了；人一旦得知天主非物质、不可见的本性，就必然以这样的前提视其为绝对不存在。反之，若这些特征在祂身上的欠缺并不构成祂存在的任何限制，那么人的心智又怎能随着物质每一属性的逐一撤除而被挤压出存在呢？\par

我反驳说：那么，我们这样辩论不过是以一个悖论换取另一个悖论；因为如果我们承认，除非撤除所有感官材料，两者都无从思考，那么我们的理性便会归结到天主与人的心智同一的结论。\par

她回答说：不要这样说；这样说也是亵渎。反之，要如圣经所告诉你的，说一个\Emph{相像}于另一个。因为那按天主的\BibleQuote{按肖像受造的}必然在各方面与它的原型相似；它在理智、非物质、与任何重量概念无关、超越任何尺寸度量方面相似；但就它自己的特殊本性而言，它与那一位不同。的确，如果它完全与那一位同一，它就不再是\Quote{肖像}了；但在那未受造的原型中我们有的是\Emph{大写的A}，在肖像中我们有的是\Emph{小写的a}；正如在一粒微小的玻璃屑中，当它恰巧面对阳光时，常可见到太阳的整个圆盘，并非按它应有的比例呈现，而是按微粒所能容许的程度呈现。同样，天主那些不可言喻的属性的反映，在我们本性的狭隘界限内闪耀；因此，我们的理性跟随这些反映的引导，通过从问题中清除所有有形性质，就不会错过把握心智的本质；另一方面，它也不会将那纯粹无限的存在降到可朽坏渺小的层次；它将视心智的这个本质仅为思维对象，因为它是这样一个存在的\Quote{肖像}；但它不会宣称这肖像与原型同一。正如我们因为宇宙中展示了神圣奥秘的智慧，而毫不怀疑有一个神圣的存有和神圣的能力存在于万有之中，维护其延续（尽管你若要求给这存有下定义，就会在其中发现天主与受造之物——无论是感官的还是思想的——完全分离，而在后者中也承认自然的区别），同样，灵魂作为实体（无论我们认为这实体是什么）的独立存在，并不妨碍它的实际存在，这一点也不奇怪，尽管世界的元素原子在其存在的定义上与它不相协调。在我们有生命的身体中，由这些原子的混合组成，正如刚才所说，在实体方面，灵魂的单纯不可见性与那些身体的粗重之间没有任何共同之处；但尽管如此，毫无疑问，身体中存在着灵魂赋予生命的影响，由一条超越人类理解力的法则所行使。甚至当那些原子再次分解为自身时，那赋予生命的纽带也并未消失；然而，正如当身体的架构仍然结合在一起时，每个单独部分都拥有一个灵魂，它平等地渗透每一个构成成员，人们不能称那灵魂是坚硬或有抵抗力的，尽管与固体混合，也不能称它为潮湿、寒冷或相反，尽管它向所有这类部分传递生命，同样，当那架构解体并回归它的同类元素时，没有什么与概率相悖，即那单纯非组合的本质，一旦通过某种不可解释的法则与身体架构一同成长，就会持续留在与它混合的原子旁边，并且不会以任何方式与已形成的结合分离。因为不能因为组合物被分解，非组合物就必然随之分解。\par

我接着说道，无人会否认这些原子会结合与分离，且此为身体之形成与消散。但需考虑：原子间有巨大间距，它们在位置、特性及特质上彼此不同。当所有这些原子汇聚于某一对象时，那理智而无维度的本质——我们称之为灵魂——与如此结合之物相合是合理的；但一旦这些原子彼此分离，各归其类，当灵魂的容器如此四散时，灵魂将如何？正如水手在船只破碎后无法同时漂浮于所有四散海面的碎片上（他只会抓住任何到手的碎片，任其余漂散），同样，灵魂本性上不能与原子一同消散，若她难以完全脱离身体，便会紧附于其中之一；若持此观点，则我们不能再将灵魂视为不朽（因她只存于一个原子），亦不能视其为可朽（因她不存于多个原子），这在逻辑上都是不一致的。\par

她答道：但理智而无维度的本质既不被压缩也不被扩散（压缩与扩散仅为身体之属性）；它因无形无体之本性，在原子聚合与分散时均同等临在于身体，既不会因原子结合之压缩而缩小，也不会因它们各自归于同类而离弃，即便异类原子间被认为存在巨大间隔。例如，轻浮与沉重、热与冷、潮湿与干燥之间差异甚大；但理智本质临在于每个曾与之结合的元素中并不费力；与对立物的混合不会使其分裂。这些元素原子在位置与特性上相距甚远；但无维度之本性依附于空间上分离之物并非难事，因为即便现在，心智能同时默观天上之物，并将其忙碌的审视延伸至地平线以外，而其默观能力并不因如此远距离的漫游而分散。因此，无论身体原子是结合还是分解，灵魂临在于其中皆无阻碍。正如金银合金中，某种有序之力熔合了金属，若随后将二者熔解分离，此种有序法则仍存于各自之中；合金分离时，此法则并不随之分裂（因不可分割者不能分成部分）。同样，灵魂的理智本质在原子聚合时可被察觉，在它们分散时亦不分裂；它与之同在，甚至在分离时亦与之同延，但自身并不被分割或按原子数目划分。此类情况属于物质与空间世界，而理智而无维度者不受空间限制。因此，灵魂存在于她曾赋予生命的原子中，没有任何力量能将她从其结合中剥离。那么，这有什么可悲伤的呢？可见之物换为不可见；为何你的理智对死亡如此憎恨？\par

对此，我回到她先前给出的灵魂定义，并说在我看来，她的定义并未足够清楚地指明灵魂所有可观察的官能。定义称灵魂是赋予有机身体生命力的理智本质，感官借此运作。然而，灵魂不仅在我们科学与思辨理智中如此运作；它不仅在那世界中产生结果，也不仅将感官用于其自然工作。相反，我们在本性中观察到许多欲望的激动和许多愤怒的激动；两者皆作为我们种类的品质存在于我们内，且在其表现中进一步显现出众多极微妙的差异。例如，许多状态由欲望引发；另有许多则源于愤怒；它们中没有一个是身体性的；而非身体者显然是理智的。我们的定义将灵魂描述为某种理智之物；因此，当我们将此观点推到结论时，必然出现两个荒谬的选项之一：要么愤怒与欲望是我们内的第二个灵魂，多个灵魂取代单一灵魂；要么我们内的思维官能不能被视为灵魂（如果\Emph{它们}不是），因为理智元素同样附于它们全部，并给它们全部打上灵魂的印记，或者同样将每一个排除在灵魂的特性之外。\par

她回答说：\Quote{你提出这个问题十分合理，此前已有许多人在别处讨论过：即我们该如何看待在我们里面的欲望原则和忿怒原则。它们是与灵魂\OriginalTerm{同质}{consubstantial}，从灵魂最初的组织就内在于灵魂自身，还是后来才附加到我们身上的不同事物？事实上，虽然所有人都同样承认这些原则可以在灵魂中察知，但研究尚未确切发现我们该如何看待它们，以便对它们获得某种固定的信念。大多数人在此问题上仍然摇摆不定，他们的意见既错误又繁多。至于我们自己，如果外邦人的哲学能系统地处理所有这些要点，确实足以证明，那么在这些思辨之外再添加一篇关于灵魂的讨论就肯定是多余的。但是，后者在灵魂问题上，按照思想者所喜欢的假想推论走了多远，我们却没有这样的自由——我指的是随心所欲地断言——；我们以圣经作为每一条教理的规则和尺度；我们必须将目光固定在圣经上，只认可那些能与圣经著作意图相协调的东西。\newline{}因此，我们必须忽略柏拉图的双轮战车和拴在其上力量不同的两匹马拉的套，以及它们的驭手——这位哲学家借此寓言式地表达这些关于灵魂的事实；我们也必须忽略他那后继者以及另一位遵循艺术规则推算或然性、并勤奋研究我们眼前这个问题的哲学家所说的一切——他宣称由于这两个原则，灵魂是会死的；我们必须忽略他们之前和之后的一切，无论他们是以散文还是以诗歌进行哲学思考；我们将采纳圣经作为我们推理的向导，圣经以公理的形式规定：灵魂中没有任何卓越不是同样属于神圣本性的属性。因为凡宣称灵魂是天主肖像的人，就断言任何与祂相异的事物都处在灵魂的界限之外；相似性不能保留在与原型不同的品质中。既然在我们所考虑的这些事物中，没有一样包含在神圣本性的概念中，那么有理由推测这类事物也与灵魂不同质。\newline{}现在，试图通过辩证法的规则和那种推导并推翻结论的科学来建立我们的教理，这是一种我们请求免于参与的讨论方式，因为它是一种证明真理的薄弱且可疑的方法。事实上，人人都清楚，那种精微的辩证法具有双向的力量，既可推翻真理，也能揭穿谬误；因此，当真理与这种技巧相伴而出现时，我们甚至开始怀疑真理本身，并认为那种精巧本身是在试图影响我们的判断、颠覆真理。另一方面，如果有人愿意接受一种赤裸裸的、非三段论形式的讨论，我们将就这些要点发言，尽可能跟随圣经传统的脉络来研究它们。\newline{}那么，我们主张什么呢？我们说，理性的动物——人——能够理解并认识，这一事实最确实地得到了我们信仰之外的人的证实；如果这一定义把忿怒、欲望以及所有这类情感视为与那本性同质，它就绝不会这样勾画我们的本性。在其它情况下，一个人不会用类的属性代替种的属性来给手头的主项下定义；因此，既然欲望原则和忿怒原则在理性本性和非理性本性中都同样可观察到，就不能正确地用这种类的属性来标志种的属性。但是，在定义本性时是多余且应予排除的东西，怎么能被当作那本性的一部分，从而用来歪曲这一定义呢？\newline{}每个本质的定义都着眼于手头主项的种的属性；凡是超出那特性之外的东西都被搁置一旁，认为与所需的定义无关。然而，毫无疑义，这些忿怒和欲望的能力被承认是一切理性和非理性本性所共有的；共有的东西与特有的东西并不等同；因此，我们绝不能把这些能力归入唯独意指人性的那些能力之中：但正如一个人可以在人身上察知感觉原则以及营养和生长原则，却并不因此而动摇对他灵魂所下的定义（因为性质A在灵魂中并不妨碍性质B也在其中），同样，当一个人在人性中察觉这些忿怒和欲望的情感时，他也不能因此正当地谴责这一定义，好像它未能充分指示人的本性。}\par

我问老师：\Quote{那么，我们对此该如何看待呢？因为我仍看不出我们怎么能恰当地否认实际上存在于我们里面的这些能力。}\par

她回答说：\Quote{你看，理性与这些情欲之间有一场战斗，要努力使灵魂摆脱它们；有些人在这场战斗中取得了成功。正如我们所知，\Person{梅瑟}就是如此，他既超越了愤怒，也超越了欲望；历史在这两方面都为他的证据：他比所有人更温良（温良表明毫无愤怒，心中完全没有怨恨），并且他也不贪图那些我们常见到一般人的欲望官能所活跃追求的东西。如果这些官能是天性，并且属于人本质的内涵，那就不可能如此。因为一个完全脱离自己天性的人不可能存在。但如果\Person{梅瑟}同时既存在又不处于这些状态，那么这些状态就不是天性，也不是天性本身。因为，如果一方面，那真正属于天性的东西是受造物的本质所在，另一方面，消除这些状态又在我们能力范围内，以至于消除它们不仅无害，反而有益于天性，那么显然这些状态应被算作外在的东西，是天性的情感，而非本质；因为本质仅仅就是其所是。至于愤怒，大多数人认为是心脏周围血液的沸腾；另一些人认为是因先前所受痛苦而急于报复；我们则认为是对激怒者施加伤害的冲动。但这些描述都与灵魂的定义不相符。再者，如果我们给欲望本身下定义，我们会称它为对缺乏之物的寻求，或对愉悦享受的渴望，或对心中所爱之物不能拥有的痛苦，或对某种无法享受的快乐的状态。这些以及类似的描述都表明了欲望，但与灵魂的定义毫无关联。对于所有其他那些我们看见与灵魂有某种关系的状况，即彼此对立的状态，如怯懦与勇敢、快乐与痛苦、恐惧与轻蔑等，也是如此；每一种似乎都与欲望或愤怒的原则相近，却又各有自己的定义以表明其独特本性。例如，勇敢和轻蔑表现出易怒冲动的一定阶段；而怯懦和恐惧所产生的性情则表现了同一冲动的减弱和削弱。痛苦又从愤怒和欲望中汲取素材。因为愤怒的无力（即不能惩罚先施痛苦者）本身成为痛苦；而获取欲望对象的绝望和心中所爱之物的缺乏则在这种阴郁状态中形成。此外，痛苦的对立面，即快乐感，也像痛苦一样，在愤怒和欲望之间分配；因为快乐是两者的主要动机。我说，所有这些状态都与灵魂有某种关系，但它们并不是灵魂，而只是像从灵魂的思维部分长出的赘疣，只因附着于灵魂而被视为其部分，却并非灵魂在其本质中所是的那个实际东西。}\par

我答复童贞女说：\Quote{然而，我们看到所有这些状态对行善之人的进步提供了不小的帮助。\Person{达尼尔}的欲望是他的荣耀；\Person{丕乃哈斯}的愤怒取悦了天主。我们也听说恐惧是智慧的开端，并从\Person{保禄}得知，\InnerQuote{合乎天意的忧苦}的目标是救恩。福音命令我们轻看危险；而\InnerQuote{不因任何惊慌而害怕}无非是对勇敢的描述，而勇敢被智慧列为善的事物之一。这一切都说明，圣经并不把这些状态视为软弱；否则软弱就不会被如此用于实践美德了。}\par

我想，\Person{教师}回答说，我自己应该为这种由于对问题的不同描述而产生的混乱负责；因为我本可以更清楚地陈述它，通过引入某种后果的顺序供我们思考，但我没有这样做。然而，现在，只要可能，我们将设计出某种这样的顺序，以便我们的论文能够按照逻辑序列前进，从而不给这类矛盾留下余地。因此，我们声明，灵魂的思辨、批判和俯察世界的官能是它凭借其自身本性特有的财产，并且灵魂由此在自己内保存了神圣恩宠的肖像；因为我们的理性猜测，神性本身，无论它在其最深的本性中是什么，正是在这些事物——普遍的监督和对善恶的批判性辨别——中彰显出来。但所有那些处于边界地带、因其特殊本性而能够倾向于两个对立面之一（它们最终被确定为善或恶取决于它们被使用的方式）的灵魂元素——例如愤怒、恐惧，以及任何其他类似的人性若被剥夺便无法被研究的灵魂情感——我们全都把这些算作来自外在的附加物，因为在作为人类原型的美善中找不到这样的特性。现在，让以下陈述仅作为一项练习（在解释中）提出。我祈祷它能免于挑剔听众的嘲笑。圣经告诉我们，天主是通过某种渐进的、有序的进展来创造人的。在宇宙的基础奠定之后，正如历史所记载的，人并没有立即出现在大地上；但走兽的创造先于他，植物又先于它们。圣经由此表明，生命力按照一种等级顺序与物质世界融合；首先，它将自身注入无感觉的本性；接着，它前进到有感觉的世界；然后，上升到理智和理性的存在。因此，既然所有现存的事物要么是有形的，要么是精神的，前者分为有生命的和无生命的。所谓有生命的，我指拥有生命；而在拥有生命的事物中，有些有感觉，其余没有感觉。再者，在有感觉的事物中，有些有理性，其余没有。既然这种有感觉的生命不可能脱离作为其主体的物质而存在，而理智的生命若不生长在有感觉中也就无法体现，因此，人的创造被记载为最后发生，好像他将每一种生命形式——植物的以及走兽中可见的——都吸收到自己内。他的营养和生长源自植物生命；因为在植物中也可以看到这样的过程，即食物由根吸收，以果实和叶子的形式散发。他的有感觉的构造来自走兽的受造物。但他的思想和理性官能是不可传递的，是我们本性中一种独特的恩赐，应单独加以考虑。然而，正如这种本性具有获取物质生存所必需之物的本能——这种本能当在我们人类中显现时，我们称之为\OriginalTerm{欲望}{Appetite}——并且我们承认这属于植物形式的生命，因为我们能在植物中也注意到它，就像许多冲动自然地活动以满足自身与其同类的营养并导致发芽，同样，走兽受造物的一切特殊状况都与灵魂的理智部分混合在一起。她继续说，愤怒属于它们；恐惧属于它们；我们内的一切其他对立活动——除了理性和思想官能之外——都属于它们。唯独那被说成是我们一切生命中被拣选之果的官能，带有神圣品格的印记。但按照我们刚刚阐明的观点，这种理性官能不可能在没有感觉的情况下存在于肉身生命中，而感觉已经存在于走兽受造物中，因此，必然地，由于这一条件，我们的灵魂与那些与之交织的其他事物有接触；这些就是我们内所有被称为\Quote{情感}的现象；它们绝非为了任何邪恶的目的而被分配给人类本性（因为如果\Emph{它们}——如此深植于我们的本性中——被发现有作恶的任何必然性，那么造物主必定是邪恶的作者），但根据我们的自由意志对它们的使用，这些灵魂的情感成为德行或邪恶的工具。它们就像铁一样，正按照工匠的意愿被塑造，接受他心中所定形的任何形状，变成一把剑或某种农具。那么，假如我们的理性——我们本性中最宝贵的部分——对这些外来的情感施行统治（正如圣经在命令人统治走兽的寓言中所宣告的），那么它们中没有哪一个会在邪恶的服役中活跃；恐惧只会在我们内产生服从，愤怒产生刚毅，懦弱产生谨慎；欲望的本能将为我们带来神圣而完美的喜乐。但如果理性放下缰绳，像被缠在战车上的御者一样被拖在后面，那么这些本能就会变成凶暴，正如我们在走兽中所看到的那样。因为既然理性并不主宰它们内所植入的自然冲动，那些更易怒的动物在愤怒的统领下互相毁灭；而庞大有力的动物并没有从自己的力量中得到好处，反而由于缺乏理性而成为有理性者的奴隶。它们对快乐的欲望的活动也不被用于任何更高的事物；在它们内观察到的任何其他本能都不会给它们自己带来任何益处。同样，在我们自己内，如果这些本能没有被理性转向正确的方向，如果我们的情感控制了我们的心智，那么人就从理性的存在变成了非理性的存在，并从像神一样的理智因这些情感的力量而堕落到走兽的水平。\par

我深受这些话感动，说道：凡认真思考的人，你的解说如此连贯流畅，虽朴素无华，却充分带有正确的印记，切中真理。对于那些只精通技术性证明方法的人，仅仅一个证明就足以说服；但至于我们，我们一致认为，有一种比任何这类人为结论更可靠的东西，那就是圣经的教导所指出的；因此我认为，除了已经说过的话之外，有必要探究这受到默感的教训是否与这一切相协调。\par

她回答说：谁能否认真理只存在于有圣经见证印记的地方呢？因此，如果需要引用福音中的话来支持我们的观点，那么研究\WorkTitle{麦子与莠子的比喻}并非不合时宜。那里的家主撒了好种子（我们显然就是那\Quote{屋子}）。但那位\Quote{敌人}，趁人睡觉的时候，将无用的种子撒在可作食物的好种子中间，把莠子种在麦子正中间。两种种子一起生长；因为撒在麦子正中间的种子不可能不与麦子一同生长。但田地的管事禁止仆人收取那无用的作物，因为它们长在相反作物的根上；免得他们连那外来的杂草一同拔掉那有营养的。现在我们认为，圣经以好种子指灵魂的相应冲动，只要它们向善培育，就必然在我们内结出德行的果实。但由于其中掺杂了对那唯独在本性上为美的真正之美的判断失误的坏种子，而后者又被随之生长的幻象所遮蔽（因为欲望的主动原则并没有朝着那自然之美——它被撒在我们内的目的——发芽增长，而是改变了方向，朝着兽性和无思想的状态移动，这种对美的错误判断将它的冲动引向这个结果；同样，愤怒的种子并没有使我们刚强勇敢，而只是武装我们去与自己的同胞争斗；爱的能力放弃了它的理智对象，完全疯狂地纵情于感官的享乐；如此类推，我们其他的情感也结出了更坏的而非更好的生长物），——因此，明智的家主让这混入他种子中的生长物留在那里，以确保我们不会因为欲望连同那无用的生长物一起被拔除而完全丧失更好的希望。如果我们的本性遭受如此的残缺，还有什么能提升我们去把握天上的喜乐呢？如果爱被拿走，我们如何与天主结合？如果愤怒被熄灭，我们有什么武器来对抗仇敌？因此，家主把这些杂种种子留在我们内，不是为了它们永远压过更宝贵的作物，而是为了使土地本身（在他的比喻中，他如此称呼心）凭借其固有的推理能力，使一种生长物枯萎，使另一种多产丰盛；但如果没有这样做，他就委托火来区分作物。因此，如果一个人适度地放纵这些情感，并且把它们掌握在自己手中而不是被它们所掌握，像国王使用臣民的众多手臂那样使用它们作为工具，那么他追求美德的努力就更容易成功。但如果他成为它们的奴隶，并且像奴隶反叛主人一样，被那些奴性的思想奴役，可耻地屈服于它们，成为他本性下等之物的猎物，那么他将被迫转向他那专横的主人命令的那些工作。既然如此，我们就不应把这些灵魂的情感——它们为拥有者所用为善或为恶——称为德行或恶习。但是，每当它们的冲动指向高尚时，它们就成为可赞美之事，如对达尼尔而言的他的欲望，对丕乃哈斯而言的他的愤怒，对那些高尚哀恸者而言的他们的悲伤。但如果它们倾向卑鄙，那么它们就是——并且被称作——恶的激情。\par

她说完这话后，让讨论停顿了片刻，我心中回顾了她所说的一切；回想起她论述中先前的那个证明线索，即灵魂在身体分解后并非不可能寄居于其原子中，我再次对她说：那么那个广为人知、名声显赫的\OriginalTerm{阴府}{Hades}在哪里呢？这个词在日常生活中，在外邦人的著作中，在我们自己的著作中都频繁出现；所有人都认为，灵魂从这里迁移到那里，如同进入一个安全的场所。你总不会把你的那些原子称为那个\OriginalTerm{阴府}{Hades}吧？\par

女导师明确地回答：显然你没有完全注意论证。在谈论灵魂从可见到不可见的迁移时，我以为关于\OriginalTerm{阴府}{Hades}的问题我已经没有遗漏什么。在我看来，无论是在外邦人的著作中，还是在神圣的著作中，这个据说灵魂所在之处的词，无非是指向那个我们无法瞥见的不可见世界的过渡。\par

我接着问：那么，为什么有些人认为阴府是一个实际的地方，它把最终离开人类生活的灵魂收归己内，将它们引向自身，作为这类本性的适当容器呢？\par

\Quote{老师回答说：\Quote{好，我们的教义绝不会因这样的假设而受损。因为如果它\Emph{是}真的，如你所说，并且天穹如此连续不断地延伸，以致它把一切事物包围在自身内，大地及其周围悬在中央，所有旋转天体的运动都围绕这个固定不变的中央，那么，我断言，绝对必然的是，无论在大地上面一侧的每个原子发生什么，在相对的一侧也会发生，因为同一个实体包裹着它的整个体积。正如当太阳在\Emph{大地}之上照耀时，阴影散布在大地的下部，因为它的球形使得它不可能同时被光线环绕，并且必然，无论太阳的光线落在球体的哪个点上，如果我们沿着一条直线直径，就会发现阴影在相对的点上，如此，连续地，在光线直线的相对端，阴影绕着那个球体移动，与太阳同步，以致大地的一半和另一半轮流处于光明和黑暗之中；同样，通过这个类比，我们有理由确信，在我们的半球中观察到的原子所遭遇的任何事，在另一个半球也会遭遇。由于原子周围的环境在大地的每一侧都是同一的，我认为既不应反驳也不应偏袒那些提出反对意见的人，他们声称我们必须认为这个或那个区域是分配给释放的灵魂的。只要这个反对意见不动摇我们的核心教义，即灵魂在肉体生命之后存在，那么在我们看来，关于地点就无需争论，因为我们坚持地点只是物体的属性，而灵魂是非物质的，其本性并不必然被拘留在任何地方。}}\par

\Quote{但我问：\Quote{但若你的对手援引宗徒的话，他说在万物复原时，所有理性的受造物都要瞻望那统御万有的那一位。在致斐理伯人的书信\BibleRef{斐 2:10}中，他提到了某些\Quote{地下}的事物：\Quote{一切膝都要跪拜}他，\Quote{包括天上的、地上的和地底下的一切}。}}\par

\Quote{老师回答说：\Quote{我们将坚持我们的教义，即使我们听到他们引用这些话。因为在死后灵魂的存在上，我们已得到对手的同意，所以我们就地点问题不作反对，如我们刚才所说。}}\par

\Quote{倘若有人问宗徒在这句话中的意思，该怎么说？你会从这段经文中去除所有地点的含义吗？}\par

\Quote{她回答说：\Quote{我不认为神圣的宗徒将理性世界划分为区域，当他命名一部分在天上、一部分在地上、一部分在地下时。理性受造物有三种状态：一种从一开始就领受了非物质的生命，我们称之为天使的；另一种与肉体结合，我们称之为人的；第三种藉死亡从肉体的纠缠中释放，存在于纯粹而简单的灵魂中。现在，我认为神圣的宗徒以其深邃的智慧看到了这一点，当他启示所有这些理性存在在未来在善业中的和谐时；他将无形体的天使世界放在\Quote{天上}，将仍与身体相连的放在\Quote{地上}，将从身体释放的放在\Quote{地下}；或者，如果还有其他属于理性世界的类别（它没有被排除）；无论有人选择称后者为\Quote{魔鬼}或\Quote{鬼神}，或任何这类东西，我们都不在乎。我们确实相信，既因为普遍的看法，更因为圣经的教导，还存在另一个世界的有灵存在，他们没有我们这样的躯体，他们与善为敌，并能损害人的性命，因自由意志从更高尚的观点堕落，并因这种对善的背叛在他们自身中体现了相反的原则；这个世界，有些人说，宗徒将其添加到了\Quote{地底下的事物}之列，在那段经文中意味着，当恶有一天在漫长的时代循环中被消灭时，善的世界之外将不剩任何东西，甚至从那些恶灵中也将和谐地升起对基督主权的宣认。如果是这样，那么没有人能强迫我们从\Quote{地底下的事物}这个表述中看到地下世界的任何地点；大气均匀地覆盖着大地的每一部分，没有任何一个角落不被这种环绕的空气所覆盖。}}\par

当她说完时，我犹豫了片刻，然后说：我尚不满足于我们一直在探究的事情；在说了这一切之后，我的心中仍有疑惑；我请求让我们讨论可以回到与之前相同的推理线索，只省略我们完全同意的那部分。我说这话，是因为我认为除了最固执的辩论者之外的所有人都已被我们的辩论充分信服，不会将灵魂在身体分解之后交付于消灭与非存在，也不会因其与原子本质不同而主张它不可能存在于宇宙中的任何地方；因为，一个理智与非物质的存有无论如何不能与这些原子相合，它（到目前为止）丝毫不被阻止存在于它们之中；而我们的这一信念基于两个事实：首先，基于灵魂在此生存在于我们的身体中，尽管它与身体根本不同；其次，基于如下事实：天主，正如我们的论证所表明的，虽明显不同于可见的质料性实体，却仍渗透于一切存有中的每一个，并藉此对整体的渗透使世界保持在存有的状态中；因此遵循这些类比，我们无需认为灵魂在从形式界进入不可见者时也不复存在。但是我坚持问，在原子整体的联合体由于它们的混合而获得了某种完全不同的形式——就是灵魂实际上已经习惯的那种形式——之后，当这种形式正如我们所料随着原子的分解而消失时，灵魂又如何跟随它们呢？既然那个熟悉的形式不再持续？\par

她等了一会儿，然后说：请允许我创造一个虚构的比喻，以说明我们面前的这个问题：尽管我所假定的可能超出了可能性的范围。那么，假设在绘画技艺中，不仅能像画家常做的那样混合相反的颜色以表现某种特定的色调，而且还能再次将这混合物分离，并恢复每种颜色的天然染料。如果白色、黑色、红色、金色或任何其他颜色，曾被混合以形成给定的色调，要再次从与他者的结合中分离出来，独自存在，我们假设我们的艺术家仍然记得那种颜色的实际本性质，并且在任何情况下都不会表现出遗忘，例如红色或黑色，如果它们在相互组合后变成完全不同的颜色，然后又各自返回其天然染料。我说，我们假设我们的艺术家记得这些颜色相互混合的方式，因此知道什么颜色与给定的颜色混合，结果是什么颜色，以及当另一种颜色从组合中被排除后，（原色）由于这种释放而恢复了其独特的颜色；并且，假设需要通过组合再次产生相同的结果，那么这个过程将更加容易，因为以前的工作中已经实践过。现在，如果理性能在这一比喻中看出某种类比，我们必须藉其光照来审视手头的问题。让灵魂代表画家的这种技艺；让自然原子代表他技艺中的颜色；让由各种染料混合而成的色调以及这些颜料返回其天然状态（我们已被允许假设），分别代表原子的汇聚与分离。那么，正如我们在比喻中假设画家的技艺告诉他每种颜色在混合后返回其适当色调时的实际染料，因此他确切地知道红色、黑色以及任何其他颜色，这些颜色通过与他种颜色的特定结合方式形成了所需的色调——这一知识包括它们混合时的外观、目前处于天然状态时的外观，以及将来如果所有颜色以类似方式再次混合时的外观——同样，我们断言，灵魂知道那些原子的天然特性，这些原子的汇聚构成了灵魂自身在其内成长的身体的框架，即使在这些原子分散之后。无论它们的天然倾向与内在的排斥力如何催促它们彼此远离，并阻止每一个与其对立者混合，灵魂仍将凭借其认知能力接近每一个，并坚持不懈地依附于熟悉的原子，直到这种分裂之后的再次以相同方式汇聚，为的是已分解的身体的重新形成，这（形成）将恰当地是并且被称为复活。\par

我打断说，你在这随口的评论中似乎对复活信仰做了极好的辩护。我认为，藉此，这一教义的反对者可能会被逐渐引导，认为原子再次合并形成与之前相同的人并非绝对不可能之事。\par

老师回答说：这很正确。因为我们可能会听到这些反对者提出下列难题。\Quote{原子相似归于相似地分解入宇宙；那么，例如，存在于某某人身上的温暖，在加入普遍温暖之后，凭借什么机制再次使自己脱离与同族的联系，以便形成那个正在被\InnerQuote{重塑}的人呢？因为如果同一个个别的粒子不返回，而只取得某种同质但不同一的东西，那么你将得到某种别的东西来代替原初的那个，这样的过程将不再是复活，而仅仅是创造一个新的人。但如果同一个人要回到自身，他必须完全同一，并在其元素的每一个原子上恢复其原始构造。}\par

为应对此种异议，我继而回答说：\Quote{如我所说，关于灵魂的上述见解将会有用；即：灵魂在分解后仍存留在她最初生长出来的那些原子中，如同一个守护者看管私有财产，当这些原子与同类原子混合时，她并不舍弃它们，反而凭借其智力的细微无处不在，精准无误地辨识它们，尽管它们极其微小；当它们与同类尘土混合时，她自身也随同属于她的原子一同扩散，在它们流回宇宙时，她毫不费力地追随其全部数量，无论自然以何种方向或方式安排它们，她始终与它们同在。一旦那全权安排者发出信号，要这些分散的原子重新结合，那么，正如从同一滑轮上垂下的各种绳索同时响应中心的拉力一样，跟随灵魂作用于各个原子的这种力量，所有这些曾经如此熟悉的原子便同时汇聚，藉着灵魂形成身体的绳索，每一个都与先前的邻居结合，拥抱旧相识。}\par

接着，导师继续说道：\Quote{也许可以恰当地将以下比喻添加到已提出的那些比喻中，以表明灵魂无需太多教导就能区分属于自己的原子与异己的原子。设想一个陶匠备有大量陶土；让这陶土数量庞大；其中一部分已塑造成型为完成的器皿，其余部分仍等待塑造；再设想这些器皿并非形状相似，而是比如一个是壶，另一个是酒坛，另一个是盘子，另一个是杯子或任何其他有用的器皿；并且，不是同一个人拥有它们全部，而是让我们想象每一个都有特殊的主人。现在，只要这些器皿没有破碎，它们的主人自然能认出它们；即使它们被打成碎片，也依然能认出；因为从那些碎片中，每个人都会知道，比如，这一块属于坛子，那一块属于杯子。如果它们再次被投入未加工的陶土中，对已加工过的部分与未加工陶土的辨别会更加准确无误。个体的人就是这样的器皿；他是从普遍质料中塑造出来的，由于原子的汇聚；他以自己独有的形态展现出与同类显著的区别；当这一形态瓦解时，曾经掌管这一特定器皿的灵魂会精确地认识它，甚至从其碎片也能认出；她也不会在与其他碎片混合时丢弃这份财产，即便它被投入原子来源的尚未成形的质料中；她始终记得自己原有的、在身体形态中紧密相合时的样子，并且在分解后，她绝不会搞错，因为残留的印记引导着她。}\par

我称赞这是巧妙揭示眼前事物自然特征的好方法；并说：\Quote{这样讲并相信它确实如此，固然很好；但假如有人引用我们主关于地狱中那些人的叙述，作为与我们的探究结果不一致的反驳，我们该如何准备应答呢？}\par

导师回答说：\Quote{那圣言叙述的说法当然是物质性的；但其中仍穿插了许多暗示，以激发有学识的探究者进行更精细的研究。我的意思是，那用深渊将善人与恶人分开、使受折磨的人渴望一滴水由指头沾着送来、使今生受苦之人安息在先祖的\InnerQuote{怀抱}中，以及讲述他们先前死亡并被葬入坟墓的那一位，正是引导明智的探究者超越字面解释。因为，那富人在地狱里抬起的是什么样的眼睛？他早已将肉身的眼睛留在坟墓里了。一个脱离肉体的灵魂如何感受火焰？他渴望用一滴水凉却的又是怎样的舌头？他已失去了血肉的舌头。用来传递这滴水的是什么样的指头？安息的\InnerQuote{怀抱}是怎样一个地方？他们两人的身体都在坟墓里，他们的灵魂是脱离肉体的，并且也不由部分构成；因此，如果我们按字面理解，就无法使叙述的框架与真理相符；只有将每个细节翻译成理念世界中的对应物，我们才能做到这一点。因此，我们必须将那深渊理解为分隔不可混淆的理念、防止它们混淆的东西，而非地上的裂口。这样的裂口，无论多大，对于脱离肉体的理智来说都不难跨越；因为理智可以瞬间到达它想去的地方。}\par

那么，我问道：\Quote{那火、那深渊和图画中的其他特征，难道不是他们所说的\Emph{是}吗？}\par

我想，她答道，那福音藉着它们每一个都意指某些关于我们灵魂问题的道理。因为当圣祖首先对那富翁说：\Quote{你一生中已享用了你的福乐}，同样地谈到那穷人，说他已尽了在分受今生苦难中的本分，然后，在那之后，又说到那\OriginalTerm{深渊}{gulf}，说是他们之间的阻隔，他显然藉这些表达暗示一个非常重要的真理；而在我看来，如下所述。人的生命曾经只有一种性质；我的意思是，它只存在于善的范畴，与恶无涉。天主的第一个诫命证明了这一点；即那赐给人无限制地享受乐园所有福乐，只禁止那善恶混合、由对立物组成之物，并将死亡作为违反那点的惩罚。但是人，因自愿冲动而自由行动，离弃了那无恶混杂的命运，并招致了那由对立物混合的命运。然而天主眷顾并没有放任我们的那种轻率而不加纠正。死亡，作为违反法律的固定惩罚，必然落在违法者身上；但天主将人的生命分为两部分，即现世生命和以后那\Quote{超脱肉身}的生命；祂给第一个部分设定了最短的时间限制，而将另一个部分延长到永恒；并因祂对人的爱，祂给人选择，要在两部分中他喜欢的那部分拥有\OriginalTerm{善}{good}或\OriginalTerm{恶}{evil}中的哪一个：或在这短暂易逝的生命中，或在那些无尽的岁月中，其界限是无限。现在这些\Quote{善}和\Quote{恶}的表达是歧义的；它们有两种用法，一种涉及理智，另一种涉及感觉；有些人将凡感觉愉快的事归类为\OriginalTerm{善}{good}；其他人确信只有理智可感知的才是\OriginalTerm{善}{good}，才配得上那名。因此，那些从未运用推理能力且从未瞥见那更好道路的人，很快就在这肉身的生命中把他们的构造所能要求的善的份额消耗在饕餮上，并且不为来世留任何东西；但那些通过谨慎和冷静的计算来节约生命力量的人，在此世受感觉痛苦的事折磨，却将他们的善保留给来世，因此他们更幸福的命运延长到与永恒生命一样长。这在我看来，就是那\Quote{深渊}；不是由大地裂开造成的，而是由今世那些导致分离成相反性格的决定造成的。一旦在今生选择了快乐，且未曾因悔改治愈其轻率的人，就把善之地置于自己无法触及之处；因为他自己挖了一个不可逾越的深渊，一种无物可突破的必然。我认为，这就是为何\Person{亚巴郎}的\OriginalTerm{怀中}{bosom}被用来指灵魂那种善的状况，圣经使忍耐的斗士在那里安息。因为据记载，这位圣祖是那时以前所有出生的人中第一个用对未来的希望交换了现在的享受；他被剥夺了他起初生命所处的所有环境，寄居在外邦人中，因而以现在的烦恼购买了未来的福乐。正如我们形象地称海洋的某段回路为\Quote{怀中}，同样，圣经似乎用\Quote{怀中}这个词来表达那些无量祝福的观念，意指一个地方，所有今世有德的航行者，当他们从今世启航后，都停泊在那祝福之深渊的无波港口中。与此同时，他们对这些祝福的否定在他们身上成为火焰，焚烧灵魂，并使灵魂渴望从那圣人富足的祝福之海洋中得到一滴的清凉；然而他们得不到。如果你还考虑\Quote{\OriginalTerm{舌头}{tongue}}、\Quote{\OriginalTerm{眼睛}{eye}}、\Quote{\OriginalTerm{手指}{finger}}以及那些没有形体的灵魂对话中出现的其他身体器官的名称，你会相信我们关于它们的猜测与我们已陈述的关于灵魂的意见一致。仔细考察那些话的意义。因为正如原子的汇聚形成整个身体的本体，同样，有理由认为相同的原因也作用于完成身体每个肢体的本体。因此，如果灵魂在身体原子重新与宇宙混合时与它们同在，它不但会认识整个质量——那曾经聚集形成整个身体的——并将与它同在，而且此外，它不会不知道每个肢体的各个材料，以至于记住原子之间的哪些划分使我们四肢完全形成。因此，假设存在于整个质量中的东西也存在于质量的每个部分中，并非不可能。那么，如果一个人想到那些原子，其中身体的每个细节潜在地存在，并猜测圣经意指\Quote{手指}、\Quote{舌头}、\Quote{眼睛}等，在分解之后只存在于灵魂的领域，就不会错失可能的真理。此外，如果每个细节将思绪从对故事的物质性接受中带走，那么，我们刚才谈到的\Quote{\OriginalTerm{地狱}{hell}}也就不能被合理认为是一个这样命名的地点；而是圣经告诉我们关于某个灵魂居留的看不见的无形境况。在这个富翁与穷人的故事中，我们也学到另一个教理，与我们先前的发现密切相关。故事使那爱感官享乐的人，当他看到自己的情况是无法逃脱的，为他在世上的亲属显出关心；\Person{亚巴郎}告诉他，那些仍在肉身中的人的生活并非没有指引，如果他们愿意，他们可以在法律和先知中找到它，但他仍继续恳求那位正义的圣祖，请求要有一个来自死者的突然而令人信服的消息传给他们。\par

我遂问：那么，这里的道理是什么呢？\par

\Quote{为什么？因为拉匝禄的灵魂正专注于他现时的福乐，丝毫不回头顾念他所留下的一切；而那个富人却仿佛用胶水一般，仍然附着于感官生活，甚至在死后也不曾摆脱它，他虽已停止生活，却仍在思想中保留着血肉（因为他在请求免去亲属的苦难时，清楚地显示他尚未摆脱肉体的感觉）。在这故事的细节中（她继续说道），我以为我们的主教导我们：那些尚在肉身中生活的人，应尽可能借着德行，在某种程度上与肉身的束缚分离并解脱，以便死后不必再需要第二次死亡来洗净因这肉身胶着而残留的污秽；一旦灵魂周围的桎梏被解开，它就能迅速而无阻碍地飞升向善，不受身体痛苦的干扰。因为如果有人彻底成为肉性的人，他的灵魂的一切运动和力量都沉浸于肉欲中，即使在脱离肉身后，他也不会与这些附着分离；正如那些长期停留在恶臭地方的人，即使进入清新空气中，也无法摆脱长久停留所沾染的不快。因此，当转入无形可见之境时，喜爱肉身的人也不可能避免拖带某些肉身的污秽；由此，他们的痛苦将加剧，因为他们的灵魂已被这样的环境所物质化。我也认为，这种观点在某种程度上与有些人所说，在坟墓周围常见逝者的幽暗幻影的说法相符。如果确实如此，那就证明那灵魂对肉身生活有过度依恋，致使它即使在被迫离开肉身后，也不愿彻底飞离，不愿接受完全转变为无形；它在肉身解体后仍停留在近旁，虽已在外，却依依不舍地盘旋于其物质所在之处，并持续徘徊。}\par

随后，我沉思了后几句的意思，说道：我认为现在你所说的和我们先前对情欲的探讨之间产生了矛盾。因为，一方面，这些在我们内里的活动（即我们之前论及的愤怒、恐惧、对快乐的欲望等等）应被视为源于我们与野兽的亲缘关系；另一方面，曾断言德行在于善用这些活动，而恶习在于滥用它们；此外，我们还讨论了其他每种情欲对德行生活的实际贡献，并发现尤其是通过欲望，我们得以更接近天主，仿佛被它的锁链从地上拉向祂——我认为（我说）那部分的讨论在某种程度上与我们现在所追求的目标相悖。\par

她问：何以如此？\par

因为，当我们在净化之后，一切非理性的本能都在我们内熄灭时，欲望的原则也将不复存在，如同其他原则一样；而一旦除去欲望，似乎对更好之道的追求也将停止，因为灵魂中不再存有任何能激发我们对善的渴望的情感。\par

对此异议，她回答说，我们这样回应。思辨和批判的能力属于灵魂肖似天主的那个部分；因为我们正是藉此领悟天主本身。那么，无论是藉由此世的先见，还是来世的净化，一旦我们的灵魂脱离与受造之物的任何情感联系，便没有任何事物能阻碍它默观美本身；因为后者本质上具有一种能力，能以某种方式吸引每一个朝向它的存在。如果灵魂净化了所有恶习，它必定会处于美的领域。天主在本质上就是美；而处于纯洁状态的灵魂将与天主有亲缘关系，并将拥抱祂如同自身。每当此时，欲望的冲动就不再需要为通往美而引路。谁在黑暗中度过时光，谁就会受到对光的欲望影响；但一旦他进入光中，享受便取代了欲望，而享受的能力使欲望变得无用和过时。因此，灵魂摆脱这类情感，转而反省自身，准确认识自己的真实本性，并在自身美的镜子和形象中看见原初的美，这对我们参与善毫无损害。因为真正的肖似天主就在于：使我们的生命在某种程度上成为至高者的副本。因为像祂那样超越一切思想、远离我们内在所观察之一切的本性，其存在方式与我们今世截然不同。人拥有一个以运动为法则的体质，被其意志冲动引向特定方向；因此，他的灵魂对于面前之物与身后之物，可以说感受不同：因为望德引领向前的运动，而当运动达到望德的目标时，记忆随之而来；如果望德将灵魂引向本质上为善的事物，那么意志活动留在记忆中的印记是光明的；但如果望德以善的幻象诱骗了灵魂，而错过了卓越的道路，那么随之而来的记忆便成为羞愧，灵魂中便产生了记忆与望德之间的内战，因为后者是意志如此糟糕的向导。羞愧所表达的正是这种心态；灵魂仿佛被结果刺痛；它对轻率尝试的悔恨如同一根鞭子，使它感到刺痛，并想借助遗忘来对抗折磨者。如今，我们的本性由于缺乏善，总是追求那仍然缺少的东西，而追求尚缺之物正是欲望的习惯，我们的体质同样显示出这点，无论是被真正的善所阻挠，还是赢得了值得赢得的善。但那超越我们对善所能形成的每一个概念、超越一切其他能力的本性，丝毫不缺乏任何可被视为善的东西，它本身就是一切善的圆满；它并非仅凭参与而处于善的领域，而是它本身就是善的实体（无论我们认为善是什么）；它既不给任何升起的望德留下空间（因为望德表现为对缺席事物的活动；但正如宗徒所问的：\BibleQuote{人已拥有的，为什么还盼望呢？}），也不需要通过记忆的活动来认识事物；因为实际看见的事物无须记忆。既然这神圣的本性超越任何特定的善，而善是爱的对象，那么当祂内视自身时，祂愿意祂所拥有的，拥有祂所愿意的，不接受任何外在之物。实际上，除了恶以外，没有任何外在之物；恶，说来奇怪，其存在在于根本不存在。因为恶的起源只是对存在者的否定，而真正存在者构成了善的实体。因此，不存在于存在者中的，必定存在于非存在者中。于是，当灵魂脱去其本性固有的各种情感，获得神圣的形式，超越欲望，进入那曾被欲望驱动而朝向的目标时，它自身内不再为望德或记忆提供栖息之所。它拥有望德的目标，而记忆则因专注于享受一切善而被逐出意识：这样，灵魂模仿了上界的生命，并符合神圣本性的独特特征；除了爱，它不再留下任何习惯，爱因自然亲和力而依附于美。因为爱就是如此：对被选之物的内在情感。当灵魂变得单纯单一，完全肖似天主，发现那真正值得热情和爱的完全单纯非物质之善时，它便依附于它，并通过爱的运动和活动与之融合，按照它不断发现和把握的那样塑造自己。通过这种对善的相似同化，灵魂成了它所参与之本性的一切，因此，由于它所参与的事物本身毫无缺乏，它自身也毫无缺乏，从而从内在驱逐了欲望的活动和习惯；因为只有找不到所缺之物时，欲望才产生。我们有天主自己的宗徒为这个教导作证，他宣布我们内在的一切其他活动（甚至为善的活动）都将被抑制和终止，唯独爱没有止境。他说：\BibleQuote{先知之恩终必停止；各种知识也终必消失}，但\BibleQuote{爱德永不止息；}这等同于说它永远如此；尽管他说信德和望德至今与爱并存，但他又将爱的期限延长到它们之后，而且有充分的理由：因为望德只在所盼望之事的享受尚未获得时才运作；同样，信德是对所盼望之事不确定时的支持；因为他这样定义——\BibleQuote{所盼望之事的实质}；但当所盼望之事真正到来时，所有其他官能都归于静止，唯独爱保持活跃，找不到后继者。因此，爱是所有卓越成就中最卓越的，是法律命令中的首要者。如果灵魂达到这个目标，它将不再需要任何其他事物；它将拥抱那万有的圆满，唯有这样才能以任何方式在自身内保存天主真正幸福的印记。因为至高者的生命就是爱，因为美对于那些认识它的人来说必然是可爱的，而天主确实认识它，因此这种认识成为爱，祂所认识的那一位在本质上是美的。这种真正的美，厌腻的傲慢无法触及；没有厌腻打断这种持续爱美的能力，天主的生命将在爱中有其活动；这生命本身就是美的，并且本质上对美有一种爱的倾向，不受任何阻碍。事实上，在美中找不到任何界限，以至于爱必须随着美的任何界限而停止。美只能由其对立面终结；但当你拥有一个像这里一样在本性上不可能变坏的善时，那善将无限地持续下去。此外，既然每个存在都能吸引其同类，而人性在某种程度上肖似天主，因为其内带有与原型的某些相似之处，灵魂便因严格的必然性而被吸引到同类的天主那里。事实上，属于天主的一切必须无论如何、不惜代价地为他保留。因此，一方面，如果灵魂没有多余的负担，也没有与身体相关的烦恼压着它，它朝向吸引它的那一位的进展是甜蜜而适宜的。但另一方面，假设它被欲向的钉子钉住，从而被束缚在物质事物的习惯上——如同地震废墟中的那些人，身体被瓦砾堆压碎；让我们想象一个比喻：这些人不仅被废墟的重量压住，而且还被废墟中发现的钉子和碎片刺穿。那么，当他们的亲属将他们从废墟中拖出以接受神圣的殡葬礼时，他们会是什么情况？全身撕裂，面目全非，以可想象的最可怕的方式，堆下的钉子因拉扯的暴力而折磨他们！——我认为，当天主因祂对人的爱，以神圣的力量将属于祂的从非理性和物质的废墟中拖出时，灵魂的处境也是如此。在我看来，天主并非出于仇恨或对邪恶生活的报复而给罪人带来那些痛苦的安排；祂只是召回并吸引那为悦乐祂而存在的一切。但祂为高尚的目的吸引灵魂到祂自己——一切幸福的泉源——时，这必然给被吸引者带来痛苦的状态。正如那些从杂质中提炼黄金的人，不仅让贱金属在火中熔化，还必须让纯金与杂质一同熔化，然后当杂质被烧尽时，金子存留；同样，当恶在炼火中被烧尽时，与恶粘连的灵魂也必然在火中，直到虚假的杂质被这火消耗并消灭。如果一种较黏的泥土厚厚地涂在绳子上，然后将绳子的一端穿过一个窄孔，然后有人在另一边用力拉绳子的那端，结果必然是一方面绳子本身服从所施加的力，另一方面涂在上面的泥土因用力拉扯而被刮掉，留在孔外，而且，这正是绳子不能顺利通过孔道的原因，并且必须被拉者强力拉扯。我想，我们可以这样想象那将自己缠绕在世俗物质情感中的灵魂的痛苦挣扎，当天主将它——祂自己的那一个——吸引到祂自己时，那已经以某种方式长入其本质的外来物质必须被用力刮除，从而引起它剧烈的难忍之痛。\par

我说：\Quote{如此看来，天主作为审判者，似乎并非首先和主要地以惩罚来折磨罪人；而是如你的论证所表明的那样，祂运作仅仅是为了将善与恶分开，并将善吸引到真福的共融中。}\par

导师说：\Quote{这正是我的意思；而且痛苦的量度将取决于每个人心中的恶有多少。因为，如果一个人像我们所假设的那样，长期停留在明知被禁止的恶中，而另一个人只犯了小罪，但在审判他们邪恶习气时，却要遭受同样程度的折磨，这是不合理的；而是根据罪恶材料的数量，那令人痛苦的火焰燃烧的时间或长或短；也就是说，只要有燃料供给它。对于积累了大量材料的人，净化之火必然非常深入；但对于火所能消耗的材料较少的地方，惩罚的剧烈程度就会减轻，因为罪恶本身减少了。无论如何，恶必须从存在中消除，这样，正如我们上面所说，绝对不存在的东西就完全停止存在。既然恶的本质不可能存在于意志之外，那么，当每一个意志都安息于天主时，恶岂不是要归于完全消灭，因为没有给它留下任何容器吗？}\par

但是我问：\Quote{如果考虑到即使受折磨一年也是极大的痛苦，从这种热切的希望中能有什么帮助呢？如果那难以忍受的痛苦延长到一个时代之久，那么，对于一个其净化与整个时代相称的人来说，后来还有什么指望中的慰藉留下呢？}\par

\Quote{唔，要么我们必须设法使灵魂完全不受任何恶的玷污；要么，如果我们的情欲本性使这全然不可能，那么我们就必须设法使我们在德性上的欠缺仅限于轻微的、易于治愈的过失。因为福音在教训中区分了欠一万塔冷通的债务人和欠五百银币、五十银币乃至一文小钱的债务人，这文小钱就是\InnerQuote{最后的一文}；它宣告天主的公义审判临到所有人，并随着债务的增加而加重所需的偿还，另一方面也不忽略最小的债务。但是福音告诉我们，这些债务的偿还并非通过退还金钱，而是将欠债的人交给刑役，直到他还清全部债务；这无非意味着用痛苦的代价来偿还那不可避免的报应，即，在他生前，当他轻率地选择纯粹的、没有与之相反的痛苦掺杂的快乐时，所应承受的痛苦份额；这样，他就除去了一切异己的、罪恶的赘疣，丢弃了任何债务的羞耻，从而能站在自由和无畏中。自由是达到一种自主的、不受管辖的状态；这是天主在起初赐予我们的，但因债务带来的羞耻感而变得模糊。一切自由在本质上都是同一的，它自然趋向于自身。由此推知，凡是自由的事物都会与相似的事物结合，并且德行是一种不受管辖的事物，也就是自由的，因而一切自由的事物都会与德行结合。再者，天主是万德的源泉。因此，那些摆脱了恶的人必将与祂结合；于是，正如宗徒所说，天主将成为\InnerQuote{万有之中的万有}；因为这句话在我看来明确证实了我们已有的观点，它的意思是天主将取代一切其他事物，并存在于一切之中。因为我们现在的生活是在多种多样的条件下进行的，我们与之有关的事物很多，例如时间、空气、地点、饮食、衣著、阳光、灯光，以及其他生活必需品，这些虽然很多，都不是天主；而我们所希望的真福状态不需要其中任何一样，但天主将成为我们的一切，并代替一切，按照那存在的每一种需要相应地分施祂自己。从圣经中也明显看到，天主对那些堪当的人成为地方、家园、衣著、食物、饮料、光明、财富、权能，以及一切可想象和可命名的、能使我们生活幸福的事物。但那成为\InnerQuote{一切}的，也将\InnerQuote{在一切之内}；在我看来，圣经在这里教导了恶的完全消灭。因为，如果天主将\InnerQuote{在一切}存在的事物中，那么恶显然就不在其中；因为如果有人假设恶那时仍存在，那么\InnerQuote{天主将\InnerQuote{在一切}}的信仰如何能保持完整呢？排除恶这一事物，就破坏了\InnerQuote{一切}一词的完整性。但那位\InnerQuote{在一切之内}的天主，永远不会存在于不存在的事物中。}\par

于是我问：\Quote{那么，对于那些面对这些灾难而心灰意冷的人，我们该说些什么呢？}\par

我们会对他们说，师傅回答道，这样。\Quote{这是愚蠢的，善人们，你们竟然为生命现实这一固定序列的锁链而烦恼抱怨；你们不知道宇宙的每个安排正朝着什么目标运行。你们不知道，万物都必须按照造物主的艺术计划，以某种规律和秩序同化于天主性。事实上，理性的受造物之所以存在，正是为了天主的恩惠的财富不致闲置。全创造智慧造了这些灵魂，这些拥有自由意志的容器，如同器皿一样，正是为了这个目的：即应当有一些能力能够接受祂的恩惠，并在恩惠的涌流中不断变得更大。参与天主的恩惠所成就的奇迹就是如此：它使接受恩惠的人变得更大、更有容量；从接受能力上，它确实使接受者获得实际体积的增长，并且永不停止扩大。恩惠之泉永不止息地涌流，而分受者的本性发现它所接受的没有任何多余或无用的东西，于是将全部流入作为自身比例的扩大，并同时变得更渴望汲取更高贵的滋养，也更有能力容纳它；两者一同成长，既是在如此丰盛的恩惠中滋养而变得更大的容量，也是随着那些增长的能力的成长而如同洪水般涌来的滋养供应。因此，这个体积很可能增长到没有任何限制可以阻止的程度，以致我们不会停止成长。面对这样的前景，你们是否因为我们的本性正在沿着为我们指定的道路向目标前进而愤怒？为什么？我们的历程无法向那里奔跑，除非那使我们沉重的东西——我是指这累赘的尘俗负担——从灵魂上被抖落；我们也不能以我们本性的相应部分居住在纯洁之中，除非我们通过更好的训练，从我们在生活中对尘俗所染上的依恋习惯中洁净自己。但如果你们对这身体有任何依恋，而且从这亲爱之物中解脱令你们痛苦，那么也不要因此绝望。你们将看到这身体的包裹物，它现在在死亡中分解，将再次从同样的原子中编织出来，但不再是这种粗糙沉重的组织结构，而是将其丝线编成更精致更空灵的东西，以至于你们不仅将拥有你们所爱之物近在身旁，而且它将以一种更明亮、更迷人的美丽归还给你们。}\par

但我现在似乎觉得，我说道，复活的教义必然要进入我们的讨论；我认为这教义乍看之下既真实又可信，正如圣经所告诉我们的；因此这不会在我们之间引起争议：但由于人类理智的软弱通过我们能理解的任何论证得到进一步加强，那么最好也不要把这部分主题留下而不加哲学探讨。让我们思考，关于它应该说些什么。\par

至于我们自身思想体系之外的 thinkers，导师继续说道，他们尽管看待事物的方式各异，却都在不同点上不同程度地接近并触及了复活的教义：他们中没有一位完全与我们吻合，但也都没有完全放弃这样的期望。有些人确实在其包罗万象中贬低人性，主张灵魂会轮流成为人和非理性之物的灵魂；它会转生于各种身体，随意从人变为飞禽、鱼类或走兽，然后又返回人类。有些人甚至将这种荒谬扩展到树木和灌木，以至于他们认为木本生命与人类相应且相似；另一些人则只持此观点——灵魂从一个男人交换到另一个男人，因此人类生命通过同一些灵魂得以持续，这些灵魂数目完全相同，一代又一代地不断出生。至于我们自己，我们坚持教会的信条，并断言只取这些思辨中足以表明沉溺其中者在一定程度上与复活教义相符的部分是明智的。例如，他们说灵魂脱离此身体后会潜入某些其他身体，这种说法与我们盼望的复兴并非完全不合。因为我们的观点认为身体，无论是现在还是将来，都是由宇宙的原子构成的，外邦人也同样持有此看法。事实上，你无法想象任何身体可以不依赖于这些原子的聚合而构成。但分歧在于：我们断言同样的身体，由同样的原子构成，再次如先前一样围绕灵魂凝聚；他们则假设灵魂降临于其他身体，不仅有理性的，还有非理性的，甚至无知觉的；并且大家都同意灵魂重新取回的身体是从宇宙的原子中获得其质料，但他们与我们分道扬镳之处在于认为这些身体不是由那些在此世生命期间围绕灵魂的完全相同的原子所构成。那么，就让这外在的见证表明灵魂再次居于身体并非不合情理；然而此后，我们有责任审视他们立场中的矛盾之处，这样通过遵循一致观点所产生的后果，就容易揭示真理。那么，关于这些理论该说什么呢？那就是：那些主张灵魂迁移到彼此不同本质中的人，在我看来，抹杀了所有自然的区别；他们在一切方面混合和混淆了理性、非理性、有感觉和无感觉之物；如果所有这些都要互相转化，而没有自然秩序的界限阻止它们相互转变。说同一个灵魂，因身体环境的特定，一时是理性而有理智的灵魂，一时又与爬虫同居洞穴，或与飞鸟结群，或成为驮兽，或成为肉食动物，或游于深海；甚或降至无知觉之物，以致生根或长成完整的树，从枝上发芽，再从芽长出花、刺、可食或有害的果实——这样说，无异于使万物相同，并相信单一的本性贯穿所有存在；认为它们之间有一种联系，彻底混合并混淆了所有可资区分的标志。断言同一事物可能生于任何事物中的哲学家，其意图无非是使万物成为一体；对他而言，事物中观察到的差异并不妨碍他将全然不相容的东西混合在一起。这就使得人们即使在看到发射毒液或肉食的生物时，也必定视其为同族，甚至同家，尽管外表不同。怀有这种信念的人，甚至会视毒参为与自己的本性不疏远，因为他会在植物中察觉到人性。葡萄串本身，虽为维持生命而栽培，他也不会毫无怀疑地看待；因为它也来自植物：我们赖以生存的麦穗果实也是植物；那么，人们怎能举起镰刀割下它们？怎能挤压葡萄串？怎能从田里拔起蓟草？怎能采集花朵？怎能捕鸟？怎能点燃火葬柴堆？因为一直不确定我们是否在粗暴对待亲属、祖先或同胞，是否通过他们某个身体来点火、调酒、备餐？认为在所有这些事物中，任何一个都有人的灵魂变成了植物或动物，而没有任何标记指明哪种植物或动物曾是人，哪种来自其他起源——这种想法会使怀有它的人对一切事物同等关心：他必须要么对活在世上的真人变得冷酷，要么如果他的天性倾向于爱亲属，他就会对每一种生命都抱同样感情，无论遇到的是爬虫还是野兽。何以，若持此论者进入一片树林，他甚至会将树木视为一群人。他的生活会是什么样？他得要么出于亲属关系而对一切温柔，要么因看不出人与其他受造物的区别而对人类冷酷。那么，根据已经说过的，我们必须拒绝这个理论；此外还有许多其他考虑，仅从一致性出发就使我们远离它。因为我曾听到持有这些观点的人说，整个民族的灵魂隐藏在某处它们自己的领域里，过着类似于具身灵魂的生活；但它们的实体如此精细轻盈，以至于它们本身随着宇宙的旋转而转动；这些灵魂因某种趋向邪恶的倾向而各自失去了翅膀，便成为具身的；首先这发生在人身上；然后，由于它们情欲的兽性亲和力，从人类生活过渡到野兽的层次；接着，又从那里降到这种无情识的纯粹自然生活，你们已经听到很多了；以至于那原本精细轻盈的灵性之物因某些恶行而变重并向下倾斜，从而迁移到人的身体；然后它的理性能力熄灭，在某种野兽中继续生活；再然后连感觉的恩赐也被收回，它变为无情识的植物生命；但之后又通过同样的等级上升，直到恢复其在天上的位置。这个教义即使经过粗略审视也会立即发现它自身没有连贯性。因为，首先，既然灵魂要因那里的邪恶从天上生活被拖到树木状态，然后又因那里所显现的德行从此处返回天上，他们的理论将无法决定哪种生活更可取，是天上的生活还是树中的生活。事实上，同一序列的循环将永远重复，灵魂无论处于哪一点，都无安息之处。如果它这样从未具身状态落到具身状态，再从那里到无情识状态，然后跳回未具身状态，那么在这些如此教导的人心中，善与恶必然产生无法解决的混乱。因为天上的生活将不再保持其福乐（既然邪恶能触及天上居民），树中的生活也不会缺乏德行（因为据说正是从这里开始灵魂向善的反弹，而从那里开始邪恶的生活）。其次，既然灵魂在天上旋转时在那里被邪恶缠绕，并因此被拖下来活在纯粹物质中，然而又从那里被提升到其高处的居所，那么这些哲学家所建立的正是他们自己观点的对立面；他们建立的是：物质中的生活是邪恶的净化，而与星辰一起的恒定旋转则是每个灵魂中邪恶的基础和原因：如果在这里灵魂借着德行生长翅膀并向上飞翔，而在那里那些翅膀因邪恶而脱落，以致它下降并附着于这个低层世界，与物质本性的粗重混杂。但这个观点的站不住脚还不止于此，即它包含了彼此截然对立的断言。除此之外，他们的基本概念本身也不能在每个方面都稳固。例如，他们说天上的本性是不变的。那么，不变中怎能容得下任何软弱？如果低层本性易于软弱，那么在这软弱之中怎能达成摆脱它？他们试图融合两个永远不能结合的东西：他们在软弱中看到力量，在情欲中看到无欲。但即使对这最后一点，他们也未始终忠实；因为他们将灵魂从物质生活带回到那个因那里的邪恶而被驱逐的地方，好像那地方的生活完全安全且未受污染；显然完全忘记了灵魂在陷入这个低层世界之前，\Emph{那里}就已经被邪恶所压。对这里下层生活的指责和对天上事物的赞扬就这样互换和颠倒；因为曾受指责的东西在他们看来通向更光明的生命，而被认为更好的状态则推动灵魂趋向邪恶的倾向。因此，要从信仰的教义中驱逐所有关于此类事情的错误和游移不定的假设！我们也不应跟随那些认为灵魂从女人的身体转移而活在男人中，或者反过来，那些离开男人身体的灵魂存在于女人中的人，仿佛他们咬住了真理；甚至如果他们只说灵魂从男人进入男人，或从女人进入女人。至于前一个理论，不仅因为它游移不定、虚幻不实，并使我们陷入彼此截然对立的观点而应被拒绝；而且它也必须被拒绝，因为它是一个不虔敬的理论，主张自然界中没有一样东西的产生不是从邪恶作为其源头获得其特殊构造的。也就是说，如果无论是人、植物还是牲畜，除非某个来自上方的灵魂落入它们之中，否则都不能出生，而这次坠落是由于某种趋向邪恶的倾向，那么他们显然认为邪恶控制着所有存在的受造。而且，以某种神秘的方式，两件事要同时发生：因婚姻而产生的人，以及与此婚姻进程同步的灵魂坠落。甚至还有比这更大的荒谬：如果，事实上，大多数牲畜在春天交配，那么，我们是否要说，春天导致在上方的旋转世界产生邪恶，以便在同一时刻，\Emph{那里}某些灵魂被邪恶浸染而坠落，而\Emph{这里}某些牲畜怀孕？关于将葡萄藤插入土中的农夫，我们又该说什么？他的手如何设法将人的灵魂与植物一起覆盖？翅膀的脱落又如何与他从事种植同时发生？同样的荒谬也存在于另一个理论中；我指的是认为灵魂必须关心已婚者的交合，并留意分娩的时机，以便它能潜入那时产生的身体。假如男人拒绝结合，或女人使自己免于成为母亲的必要，那么邪恶是否会无法压住那个特定的灵魂？婚姻是否会在上方奏响灵魂中邪恶的第一个音符？还是这种颠倒的状态会完全独立于任何婚姻而侵袭灵魂？但那么，在这种情况下，灵魂将不得不像个无家可归的流浪汉一样在间隔中游荡，既已从天上的环境坠落，又可能在某些情况下仍无身体接纳它。此后，他们怎能想象天主的眷顾遍及世界呢？因为他们将人类生命的开始归因于这种偶然且无意义的灵魂坠落。因为所有后续必然与开始相符；所以，如果生命因偶然事故开始，其整个进程立即成为一连串事故，而试图使整个世界依赖于神圣力量，对这些否认其中个体由神圣意志之命而生、并将各个存在的起源归因于来自邪恶的遭遇的人而言，是荒谬的，仿佛除非邪恶敲响了它的基调，否则就不可能有人的生命。如果开始是这样，那么续篇必定会根据那个开始而展开。没有人敢声称美丽可以从肮脏中产生，正如善不能产生其反面。我们期望果实与种子的本性相符。因此，这盲目的偶然运动将统治整个生命，而没有任何天主的眷顾再遍及世界。\par

不但如此，甚至我们计算的预测也将全然无用；德行将失去其价值；离弃邪恶也将不值得。\newline{}一切将完全受制于掌舵者机缘；我们的生命将毫不异于没有压舱物的船，在无法预测的境遇的波浪上漂流，时而碰上这件好事，时而碰上那件恶事。\newline{}德行的宝藏绝不会在那些其体质归因于与德行完全相反的原因的人身上找到。\newline{}如果天主确实统管我们的生命，那么，显然地，邪恶不能起始生命。但如果我们的出生确实归因于邪恶，那么我们就必须以完全与之相称的方式继续生活。\newline{}由此将显示，谈论死后等待我们的\Quote{改造之所}和\Quote{公正的报应}，以及所有其他那里所主张并相信的有助于压制恶习的事，是愚蠢的；因为一个人既然其出生归因于邪恶，又怎能置身其外？一个人，如其本质起源于一种恶习（正如他们所断言），又怎能拥有任何朝向德行生活的自觉冲动？\newline{}拿任何一个受造的无理性动物来说；它并不尝试像人一样说话，而是使用似乎是随着母乳吸入的那种天然发声方式，并认为失去清晰的言语并非损失。\newline{}同样地，那些相信一种恶习是他们生命的起源和原因的人，永远不会使自己渴望德行，因为德行是完全悖离他们本性的事。\newline{}但事实上，那些藉着反思洁净了灵魂视觉的人，都渴望并追求德行生活。因此，这一事实清楚地证明，恶习在时间上并不先于生存的开端，我们的本性也并非源于恶习，而是天主统管万有的智慧是其原因：简而言之，灵魂以令其造物主喜悦的方式登上生命的舞台，然后（而非之前）藉其自由意志，自由选择它心中所想之事，因此，无论它愿意成为什么，就成为什么。\newline{}我们可以藉着眼睛的例子理解这一真理。看是它们的自然状态；但看不见要么出于选择，要么出于疾病。这种非自然状态可以取代自然状态，要么是故意闭眼，要么是因疾病失去视力。\newline{}同样地，我们可以断言，灵魂的构成源于天主，既然我们不能想象天主有任何恶习，灵魂便免于任何必然邪恶；然而，尽管灵魂是在这种状态中产生的，它却可以凭自己的自由意志被吸引向一个选定的方向，要么故意对善闭眼，要么让它们被那个我们引以为伴的阴险敌人损伤，从而在错误的黑暗中度过一生；或者相反，保持对真理的视线清晰，并远离一切可能使它暗淡的软弱。——但有人会问：\Quote{它何时及如何产生的？}\newline{}至于某一事物如何产生的问题，我们必须完全将其排除在讨论之外。即使对于那些完全在我们理解范围内且我们有感性知觉的事物，思辨理性也不可能把握现象产生的\Quote{如何}；以至于即使是受启示的圣人也认为此类问题无法解决。例如，宗徒说：\BibleQuote{因着信德，我们明白诸世界是藉天主的话造成的，这样，看得见的是从看不见的造出来的\BibleRef{希 11:3}。}他绝不会，我认为，如果他以为这个问题可以藉理性能力的任何努力来解决，就不会那样说了。\newline{}宗徒一方面确认，世界本身及其中万物是藉天主的旨意造成的（无论这\Quote{世界}如何包含整个有形与无形受造界的概念），另一方面他却将这一造法的\Quote{如何}排除在探究之外。我也不认为任何探究者能够触及这一点。这个问题表面呈现出许多无法克服的困难。\newline{}例如，一个运动的世界怎能来自一个静止的世界？有维度与复合性的怎能来自单纯与无维度的？它确实来自至高存有吗？但这个世界与那存有种类不同的事实，反对这样的假设。那么它是来自其他源头吗？然而，信德无法设想任何完全在神圣本性之外的事物；因为在创造因之外，我们若假设某种另物，工匠以其全部技巧必须借助它以形成宇宙，那么我们就必须相信两个不同且分离的本原。\newline{}既然万有的原因是唯一且独一的，而由这原因所产生的存在却不与其超然特质的本性相同，那么在我们的两种假设中便出现了同样巨大的不可思议性，\Emph{即}：创造直接出自神圣存有，以及宇宙的存有归因于祂以外的某个原因；因为如果受造物与天主具有相同本性，我们就必须认为祂具备属于其受造物的属性；否则，一个物质世界，在天主实体的范围之外，且因其缺乏一切开端而与自有者的永恒平等，就必须被列为与祂对立：而这正是摩尼的追随者以及某些持同样大胆见解的希腊哲学家所想象的；他们并将这一想象提升为一个体系。\newline{}那么，为了避免陷入探究事物起源所涉及的这两种荒谬之中，让我们效法宗徒的榜样，对每个受造物中的\Quote{如何}问题完全不干预，而仅仅顺便观察：天主圣意的运动在任何祂愿意的时刻成为事实，且其意图立刻在受造界中实现；因为全能者不会将其远见之智的计划停留在无形愿望的状态：愿望的实现就是实体。\newline{}简言之，整个存在物世界分为两类：\Emph{即}理性受造界和物质受造界；而理性受造界起初似乎丝毫也不与精神存有相冲突，反而与之相近，因为它表现出我们正确归之于祂超然本性的那缺乏可触形式和维度。\newline{}另一方面，物质受造界必然被归类为与神性毫无共同之处的特殊物；并且它确实给理性提出了最大的困难，即：理性无法看见有形者\Emph{如何}出自无形者，坚硬固体\Emph{如何}出自不可触者，有限\Emph{如何}出自无限，受一定比例限制、引入了数量概念的事物\Emph{如何}出自无大小、无比例者，等等，贯穿身体的每一个具体情况。\newline{}但即使对于这一点，我们也可以说这么多：\Emph{即}，我们归于身体的那些事物中没有一样是身体本身；无论是形状、颜色、重量、广延、数量，还是任何其他限定概念；它们每一个都是一个范畴；将所有这些结合成一个单一整体才构成身体。那么，既然这些完成特定身体的各个限定仅由思想把握，而非感觉，并且神性是一个思想者，那么对于这样一个思想主体来说，产生那些相互结合为我们生成身体实体的可思考之物，又有什么困难呢？\newline{}然而，所有这些讨论都超出了我们当前的事务。前一个问题是一如果有些灵魂存在于其身体之前，那么它们何时及如何产生？而对于这个问题，关于\Emph{如何}的部分，我们已经排除在考察之外，并未干预，因为它呈现不可逾越的困难。剩下的问题是灵魂开始存在的\Emph{何时}：它紧接在我们已经讨论的之后。因为如果我们承认灵魂在身体之前已存在于其特有的某处居所，那么我们就无法避免看到最近讨论的那些奇异教导中有某种力量，这些教导将灵魂居住在身体中解释为某种恶习的结果。\newline{}再者，另一方面，任何能反思的人都不会想象灵魂的出生晚于身体的形成，即灵魂比身体的形成更迟；因为每个人都能亲眼看到，所有无生命或无灵魂的事物中没有一个拥有任何运动或生长的能力；然而，在子宫中孕育的东西既能生长又能在原地移动，这是没有疑问的。因此，剩下的就是我们必然认为身体和灵魂的生存起点是同一的。\newline{}我们也肯定，正如大地从农人手中接收树苗使之成树，它自身并不赋予其幼苗生长的能力，而只是在它被放入其中时借给它生长的冲动，同样地，从人分泌出来用以播植人的物质本身在某种程度上是一个有生命的存有，与其来源一样赋有灵魂并能滋养自身。\newline{}如果这个幼苗因其微小而不能一开始就包含灵魂的所有活动与运动，我们也无需惊讶；因为在谷物的种子里也并非立刻可见穗。如此大的东西怎能挤进如此小的空间呢？\newline{}但大地不断用适宜的养料喂养它，于是谷粒便成为穗，在土块中并未改变其本性，而是在养料的刺激下发展并使之完美。\newline{}因此，正如那些生长中的种子向完美的进展是渐进的，同样，在人的形成中，灵魂的力量按照其身体所达到的大小而显示自身。它们首先在胎儿中显现，表现为营养和发展的能力；此后，它们将知觉的恩赐引入已出生的有机体；然后，当这达到后，它们显露出一定程度的理性官能，如同某种成熟植物的果实，并非一下子全部生长，而是随着该植物的抽芽而持续进步。\newline{}由此看到，从一个有生命者分泌出来为奠定另一个有生命者基础的东西本身不可能是死的（因为死亡状态源于生命的缺乏，而缺乏不可能先于拥有），我们从这些思考中把握到这一事实：在由两者（灵魂与身体）结合而成的复合体中，两者同时进入存在；一个既不先来，另一个也不后到。\newline{}但关于灵魂的数目，我们的理性必须考虑其增长终将停止；这样，自然之流就不会永远流淌，在连续生育中向前倾泻而永不停止其前进运动。我们种族有朝一日会停下来的原因，依我们的看法，如下：既然每一个理智性的实在都固定于自身的圆满中，那么期望人类也将达到一个目的是合理的（因为在这方面人类也不与理智世界分离）；因此我们应当相信，人类不会永远像现在这样仅在不完善中可见：因为后代的持续增加表明我们种族中存在某种匮乏。\par

如此，当人类抵达它固有的完满时，这种生生不息的涌流运动将会完全停止；它将触及命定的界限，一种与现今生死更替截然不同的新秩序将延续人类生命。若无出生，则必然无死亡。分解必先有组合（我所谓组合即藉出生而进入此世）；因此，若这合成不曾先存，便无分解可随。因此，若依或然之理，在此世之后的生命向我们预示为固定而不朽的，没有出生与朽坏来改变它。\par

导师结束了她的阐述；坐在她床边的许多人看来，整场讨论已得出了恰当的结论。然而，我担心导师的病况若发生致命转变（事实正是如此），我们之中将无人能回答不信者对复活提出的异议，于是仍坚持追问。\par

辩论尚未触及与我们信仰有关的所有问题中最关键的一点。我的意思是，受默感的著作（包括新约和旧约）不仅极其强调地宣告：当我们的种族藉着世代更替、在世代完成其完整循环后，这生生不息的洪流将完全停止；而且，当圆满的宇宙不再容许增长时，所有灵魂将以其全部数目，从不可见且分散的状态回归到可感知与光明之中，属于每个灵魂的同一原子依原先的次序重聚；这种人类生命的重建，在这些包含天主教导的著作中称为复活，原子整个运动的名称与从地上实际扶起者的名称相同。\par

但她说：\Quote{在已说过的话中，哪一点被忽略了呢？}\par

我答道：\Quote{嗯，正是复活的切实教义。}\par

她回答说：\Quote{然而，在我们漫长而详细的讨论中，许多地方都指向了这一点。}\par

我坚持说：\Quote{那么，难道你不知道我们的对手针对你那望德所提出的所有异议，如蜂群一般繁多吗？}\par

于是我立刻试着复述那些好辩者为了推翻复活教义而想出的全部伎俩。\par

她回答说：\Quote{我认为，首先，我们必须简要浏览圣经中有关这一教义的零散宣告；它们将为我们的论述画上句号。}\newline{}那么，请注意，我听见\Person{达味}在神圣歌咏的赞美中，在\BibleBook{}{圣咏}第103（104）篇的赞歌结束时说，他以上主对世界的治理为主题，\BibleQuote{祢收回它们的气息，它们就死亡，归于尘土；祢派遣祢的圣神，它们就被创造；祢使地面更新。}\newline{}他说，那在万物中运行的圣神的能力，使它所进入的存在获得生命，而它离弃的存在则失去生命。\newline{}既然经上宣告死亡发生于圣神的离去，而这些死者的更新发生于圣神的显现，而且按陈述的顺序，那些将要如此更新的人的死先出现，我们便认为，在这些话中，复活的奥迹被宣告给了教会，而\Person{达味}以预言的精神表达了你们正在询问的这一恩赐。\newline{}你会在另一个地方发现同一位先知也说：\BibleQuote{世界之天主、万有之主已向我们显现，好让我们在装饰者中间过节}；他提到用树枝的\Quote{装饰}，指的是搭帐篷节，按\Person{梅瑟}的诫命，这节日自古就遵守。\newline{}我认为，那位立法者采用了先知的精神，在其中预言了未来的事；因为尽管装饰一直在进行，却从未完成。\newline{}真理确实在那些不断发生的节期的预象和谜语中被预示，但真正的搭帐篷尚未到来；因此，按先知的宣告，\BibleQuote{整个世界之天主和主已向我们显现，好让我们这个已解散的住所的搭帐篷得以为人类举行}；也就是说，一个物质的装饰可能通过我们分散的原子聚合而重新开始。\newline{}因为那个词\OriginalTerm{装饰}{\GreekText{πυκασμὸς}}在其特殊含义上指圣殿的围场和完成它的装饰。\newline{}现在\BibleBook{}{圣咏}的段落如下：\BibleQuote{天主和主已向我们显现；在装饰者中间过节，直到祭台的角}；在我看来，这以隐喻宣告了一个事实：整个理性受造物要守同一个节期，而且在那圣者的集会中，低微者将与他们之上的者一同舞蹈。\newline{}因为在那作为预象的圣殿的构造中，并非所有在它围墙之外的人都允许进入，而一切外邦人和外邦人皆被禁止进入；此外，在那些已进入的人中，并非所有人都同样有特权前进到中心；只有那些以更圣洁的生活方式和某些洒净而奉献自己的人才能；而且，在这些最后的人中，并非人人都能踏足圣殿的内部；只有司铎有权进入帷幔内，且仅为圣所服务；而即使是司铎，那昏暗的圣殿内室——那里有美丽的祭台及其突出的角——也是禁止进入的，除了他们中的一位，他担任司铎的最高职务，每年一次，在一指定的日子，独自进入，携带比通常更神圣和奥妙的祭品。\newline{}既然你们所知道的这座圣殿有这些区别，它显然是精神世界状况的表象和模仿，这些物质礼规所教导的教训是：并非整个理性受造物都能接近天主的圣殿，换言之，全能者的朝拜；相反，那些被虚假信念引入歧途的在天主之域之外；而在那些因朝拜而被拣选并得以进入的人中，有些人因洒净和净化而享有更多特权；而在这些最后的人中，那些被祝圣为司铎的人有更进一步的特权，甚至被允许进入内室的奥迹。\newline{}而且，为使寓言的含意更加清楚，我们可以理解圣言在此教导：在一切赋有理性的能力中，有些像圣祭台一样固定在天主的最内层圣所；而这些最后的能力中有些像角一样突出，因其卓越；在它们周围，其它的能力按照规定的等级次序排列为第一或第二；相反，人类种族，由于内在的邪恶，被排除在天主之域之外，但用净化水净化后，又重新进入；并且，既然我们罪孽所筑起的一切进一步屏障将最终被拆除，当那时刻来临，我们本性的帐篷仿佛要在复活中重新搭建，而一切罪恶的积习从世界消失时，那时复活中已装饰自己的人将在天主周围举行普世节庆；同一宴席将为所有人摆设，没有分别使任何理性受造物无法平等参与；因为那些现今因罪而被排除的人，最终将被接纳进入天主福乐的最神圣之处，并将自己系于那里祭台的角，即那些超越性能力中最卓越者。\newline{}宗徒更明确地说出了同一件事，当他指出整个宇宙最终与善和谐一致时：\BibleQuote{为使一切在天上的、地上的和地下的，都向祂屈膝，且一切唇舌都宣认耶稣基督是主，以光荣天主圣父}；他用\Quote{角}指那些属天使的、\Quote{在天上的}事物，用其他术语指我们自身，即我们接下来想到的受造物；一个联合之声的节庆将占据我们所有人；那节庆将是对那真正存在者的宣认和认识。\newline{}她继续说道：\Quote{人们可以选取圣经的许多其他段落来确立复活的教义。}\newline{}例如，\Person{厄则克耳}以预言的精神跃过所有中间的时间及其巨大跨度；他凭预见的能书立在实际复活的时刻，仿佛他真的看见了未来之事，将其描述呈现在我们眼前。\newline{}他看见一片广阔平原\BibleRef{则 37:1}，展开至无尽远处，平原上散落着大堆骨头，东一堆西一堆；然后在天主的推动下，这些骨头开始移动，与它们曾有的同类聚在一起，附着于熟悉的关节，然后用肌肉、肉和皮肤包裹自己（这就是\BibleBook{}{圣咏}诗歌中称为\Quote{装饰}的过程）；事实上，圣神赋予那里的一切生命和活动。\newline{}但至于我们的宗徒对复活奇迹的描述，既然容易找到并阅读，何必重复？例如，用圣言的语言，\BibleQuote{在一声呼喊与号角声中}，一切死去的和倒下的都将\BibleQuote{在一眨眼间}被改变成不朽的存在。\newline{}福音中的表达我也略过；因为其含义对每个人都十分清楚；我们的主不仅在言语上宣告死者的身体将要复活；而且以行动显示复活本身，从我们更容易触及、更不易怀疑的事情开始这一奇妙工程。\newline{}首先，祂在致命的疾病中显示祂赋予生命的能力，用一句话驱散那些病患；然后祂使一个刚死的小女孩复活；然后祂使一个已被抬出的年轻人坐起来在担架上，并交给他母亲；之后祂从坟墓中唤出已死四天、已经腐烂的\Person{拉匝禄}，用祂命令性的声音使倒下的身体活起来；然后三天后，祂从死者中复活祂自己的人性身体，尽管它被钉子和长矛刺穿，祂把那钉痕和枪伤带来作为复活的见证。\newline{}但我认为详细提及这些事情并无必要；因为那些接受了这些文字记载的人心中对此毫无怀疑。\par

我说道：\Quote{但那并非问题的重点。你的大多数听众都会同意将来有一天会有复活，并且人将被带到不朽的审判台前；这既由于圣经的证明，也由于我们之前对问题的探讨。但问题依然存在：我们所期待的状态是否与身体目前的状态相似？因为如果是这样，那么，正如我所说过的，人最好根本不要期盼复活。因为如果我们的身体将要以它们停止呼吸时相同的那种状态重新获得生命，那么人在复活中可期盼的将是无尽的灾难。因为有什么景象比在极度的老年时我们的身体枯萎、变成可憎而丑陋的东西更为可怜呢？那时，肉身经年消耗殆尽，皮肤干枯贴骨，满是皱纹，肌肉因不再得到天然水分滋润而痉挛，整个身体因而收缩，两侧双手无力从事天然工作，不由自主地颤抖？长期痨病者的身体又是怎样的景象！它们与光骨头仅有的区别是给人一种披着破旧皮肤面纱的外观。水肿病肿胀者的身体又是怎样的景象！什么言语能描述麻风病患者的丑陋残损？渐渐地，腐烂蔓延并吞噬他们所有的肢体和感觉器官。什么言语能描述那些在地震、战斗或任何其他灾祸中致残，并在这种境况中苟延残喘许久才自然死亡的人？或者那些因受伤而从婴儿期起就肢体畸形的人！对于他们，人能说什么呢？人应如何看待那些曾被弃、被勒死或自然死亡的新生婴儿的身体，如果它们要被复活成原本的样子？它们会继续那种婴儿状态吗？还有什么境况比那更悲惨？或者它们会达到年富力强的状态？那么，大自然又有什么样的乳汁来再次哺乳它们呢？那么问题归结于：如果我们的身体要在每个方面都与先前相同地复活，那么我们所期盼的这件事简直就是一场灾难；而如果它们不相同，那么复活起来的人就会是不同于死去的人的另一位。例如，如果一个小男孩被埋葬，却有一个成年男子复活，或者相反，我们怎能说死者以其自身被复活了呢？因为由于年龄的差异，已有另一个人代替了他。人们看到的不是孩子，而是成年男子。不是老人，而是正值壮年的人。事实上，完全不是同一个人，而是另一个人。残疾人变成健全人；肺痨患者变成肌肉结实的人；其他所有可能的情况也是如此，为避免冗长就不一一列举了。那么，如果身体不会以其与泥土混合时的属性复活，那么那具死尸就不会复活；相反，泥土将被塑造成另一个人。那么，当另一个人而不是我复活时，复活又如何影响我呢？我说的是另一个人；因为当我看到不是我自己的人取代了曾经的我时，我怎能认出自己呢？它不可能是真正的我，除非它在每一个方面都与我相同。例如，假设今生我记忆中某人的特征：秃头，厚嘴唇，鼻子稍扁，肤色白皙，灰眼睛，白发，皮肤起皱；然后我去找这样一个人，却遇到一个头发浓密、鹰钩鼻、肤色黝黑、在其他方面容貌类型完全不同的年轻人；我看到后者时，会想起前者吗？但为什么更久地停留在这些对复活较不有力的反对意见上，而忽略了最强有力的一点呢？因为谁没有听说过，人的生命像一条河流，以一定的进度从出生流向死亡，只有当它也不再存在时才停止这种前进的运动？这种运动确实不是空间上的变化；我们的体积从不超出自身；但它通过内在的变易而前进；只要这种变易正如其名所示，它就永远不会（从一刻到下一刻）保持在同阶段；因为正在改变的事物如何能保持任何同一性呢？灯芯上的火焰，就外观而言，似乎总是相同的，其运动的连续性使它看起来像一个不间断且自足的整体；但实际上它总是在传递自身，从不保持相同；被热力提取出来的水分一燃烧起来就烧尽并变成烟，这种变易力推动火焰的运动，自行将质料变成烟；因此，正如一个人在同一地方两次触摸那火焰，却不可能两次触摸到完全相同的火焰（因为变易的速度太快，它不会等待第二次触摸，无论触摸多么迅速；火焰总是新鲜而崭新；它总是在产生，总是在传递自身，从不停留在同一位置），我们发现我们身体的结构也是同样的情况。在它内部，有流入与流出，以变易的进程进行，直到它停止存活；只要它活着，就没有停止；因为它要么被补充，要么以蒸气排出，要么被这两种过程的结合推动着。那么，如果一个人甚至与昨天的自己也不相同，而是因这种转化而变得不同，那么当复活使我们的身体重获生命时，那个单独的人将变成一群人类，以致在他的复活中，将同时出现婴儿、孩童、少年、青年、男子、父亲、老人，以及他曾经所是的一切中间阶段的人。此外，贞洁与放荡都是在肉身上发生的；那些为了信仰而忍受最痛苦折磨的人，以及那些回避这些的人，这两类人都与肉身感觉有关而显露出他们的品性；那么，在审判时怎能实行正义呢？}\par

试想同一个人先是犯罪，然后藉悔改洁净自己，以后又可能再陷于罪中；在这种情况下，玷污的身体和未玷污的身体都随着他本性的改变而改变，二者都不能终始如一。那么，放荡者将在哪个身体中受刑罚？是在那个因年老而僵硬、临近死亡的身体吗？但那个身体与犯罪的身体并非同一个。那么，是在那个因顺从情欲而玷污了自己的身体吗？但那个老年身体又在哪里？实际上，后者不会复活，复活将不会完成全部工作；否则他将会复活，而罪犯却得以逃脱。让我再从不信者对此教义所提出的反对意见中举出一点。他们主张：身体没有一部分是天性无用的。例如，有些部分是我们赖以生存的动力因；没有它们，我们在血肉中的生命就无法延续；这些部分如心、肝、脑、肺、胃及其他生命器官；另一些部分被指定用于感觉活动；另一些用于抓握和行走；还有些部分适合于延续后代。如果未来的生命与现世完全相同，那么我们所设想的改变就化为乌有；但如果那报告是真的（它确实是），即婚姻不属于那未来生命的秩序，饮食也不再维持那生命的延续，那么当我们不再期望在那生命中从事任何现有的器官所服务的活动时，我们的身体器官还有什么用？如果为了婚姻，现在有某些器官适于婚姻，那么当婚姻不再存在时，我们就不会需要那些器官；同样，用于工作的手、用于奔跑的脚、用于摄取食物的口、用于咀嚼的牙齿、用于消化的胃器官、用于排泄多余物的管道，也是如此。因此，当所有这些功能都不再存在时，它们的工具如何、为何存在？所以，必然地，如果那些不会以任何方式有助于那未来生命的东西不应存在于身体周围，那么构成现今身体的所有部分也将不存在。那么，那生命将借助其他工具来进行；没有人能称这种状态为复活，因为特定的肢体因对那生命无用而不再存在于身体中。但另一方面，如果我们的复活将体现在这些每一个肢体中，那么复活的创始者将在我们身上塑造那些对那生命无用且无益的东西。然而，我们不但必须相信有复活，而且必须相信复活不会是荒谬的。因此，我们必须留意对此的解释，以便这真理的每一部分都能使其合理性得以保全。\par

当我结束后，老师便回答说：你颇有精神地攻击了与复活有关的教义，以所谓的修辞学方式；你用似是而非的颠覆性论点绕着真理打转，以至于那些没有非常仔细思考这一奥秘真理的人可能会因为这些论点的似真性而受影响，并可能认为针对前述内容提出的困难并非完全无关紧要。但是，她继续说道，真理并不在于这些论点，即使我们可能无法用同样有力的语言给予修辞性的回答。所有这些问题真正解释仍然储存在智慧隐藏的宝库中，直到我们通过复活的实际经验被教导复活的奥秘时才会显明；那时，将不再需要言语来解释我们现在所希望的事物。正如在夜晚熬夜的人们可能会就阳光是什么样的东西展开许多争论，然后阳光简单而美丽的显现会使任何言语描述成为多余，同样，试图推测地达成未来状态的每一种计算，在对我们希望的客体降临时，都将化为乌有。但既然我们有责任不放过任何未加检验的反对论点，我们将按以下方式阐述关于这些要点的真理。首先，让我们清楚理解这教义的范围；换句话说，就是圣经颁布它并创造对它的信仰的最终目的是什么。那么，为了勾勒如此庞大真理的轮廓并用定义来把握它，我们将说：复活是\Quote{\Emph{我们本性在其原始形式下的重建}}。但是，在那种生命形式中，其创造者是天主自己，我们有理由相信既没有老年，也没有婴儿期，也没有任何源自我们现在各种软弱痛苦的苦难，也没有任何形式的身体折磨。我说，我们有理由相信，天主不是这一切的创造者，而是在人的人性遭受邪恶攻击之前，人是神圣的；然后，随着邪恶的入侵，所有这些苦难也侵入了他。因此，摆脱邪恶的生命完全无需在源于邪恶的事物中度过。因此，当人穿过冰层时，他的身体必然变冷；或者当他在烈日下行走时，皮肤必然变黑；但如果他避开了其中之一或两者，他就完全避免了这些结果——既不变黑也不变冷；事实上，当特定原因被移除时，没有人有理由期待特定原因的结果。同样，我们的本性变得有激情，就不得不面对激情生活的一切必然结果；但当它回归到那种无激情的福乐状态时，它将不再遭遇邪恶倾向的不可避免的结果。因此，既然所有动物生命注入我们本性的事物在我们的人性因接触邪恶而堕入激情之前并不存在，那么当我们放弃这些激情时，我们肯定会放弃它们所有可见的后果。因此，没有人有理由在另一种生命中寻找任何激情在我们身上留下的后果。正如一个人，穿着破烂的短上衣，脱下这件衣服后，就不再感到它带来的耻辱，同样，当我们脱下那件由兽皮制成的死气沉沉的丑陋短上衣（因为我认为 \Quote{皮衣} 是指当我们习惯于激情放纵时被穿上的那种属于兽性本性的构造）时，我们将在脱下那件短上衣的同时，扔掉我们周围那层兽皮的所有附属物；这些附属物包括性交、受孕、分娩、不洁、哺乳、进食、排泄、逐渐长到成熟、壮年、老年、疾病和死亡。如果那层皮不再在我们周围，它的后果怎么能留在我们里面呢？因此，当我们期待来世有不同的事态时，以与此无关的事情为由反对复活的教义是愚蠢的。我的意思是，消瘦或肥胖，消耗状态或多血症，或任何其他在流动不息的本性中出现的状态，与来世有什么关系？来世是陌生的，不像这样短暂易逝的状态。对于复活的操作，只需要一件事：即一个人必须通过出生而活过；或者用福音的话说，\BibleQuote{一个人应生在世上}；生命的长短，死亡的方式如何，是与该操作无关的问题。无论我们取什么例子，无论我们如何假设此事，都是相同的；这些生命中的差异没有给复活带来任何困难，也没有带来任何便利。一旦开始生活的人，在死亡中解体的间隙在复活中被修复之后，必然继续存在曾经活过的事实。\par

\Emph{如何}与\Emph{何时}，\Emph{它们}与复活有何相干呢？对这类问题的思考，完全属于另一条探究路线。例如，一个人可能生活在身体的安逸中，或困苦中，有德行或邪恶，有名望或羞辱；他可能悲惨或幸福地度过了一生。这些以及诸如此类的结果，必须从他的寿命长短和生活方式中得出；而要对他一生所行之事作出判断，法官必须审察他的放纵（视情况而定）、他的损失、他的疾病、他的老年、他的壮年、他的青年、他的财富、他的贫穷：一个人处于这些境遇中，究竟如何结束了他命定的生涯；他是否在漫长的一生中领受了许多祝福，或许多灾祸；或者对两者全然未尝，而在心智能力形成之前就停止了生命。但无论何时，当天主将我们的本性恢复到人的原始状态时，那时谈论这些事情，以及想象基于这些事情的异议能证明天主的能力在达到其目的时受到阻碍，将是徒劳无益的。祂的目的只有一个，唯一的一个；那就是：当我们的全人类从第一个人到最后一个人都得以完成时——有些人在今生立即被净化了邪恶，另一些人在随后的必要时期被火治愈，还有一些人在此生同样对善与恶毫无知觉——祂要赐给我们每一个人分享祂内之祝福，圣经告诉我们：\BibleQuote{眼所未见，耳所未闻}，人心从未想到过。但这并非别的，至少照我理解，就是在于天主自身之内；因为那超越听觉、视觉和心灵的善，必定是那超越宇宙的善。但现世有德行与邪恶生活的区别，将以这种方式显明：即每个人在参与那应许的福乐时，或快或慢。根据每个人根深蒂固的邪恶程度，其治疗期限将被计算。这治疗在于其灵魂的洁净，而这不能在没有一种剧痛状态的情况下实现，正如我们之前的讨论中所阐述的。但任何人若试图探究我们宗徒智慧的深度，就会更充分地明白所有这些异议的徒劳和不相干。当他向格林多人解释这个奥迹时（格林多人可能自己提出了与今天攻击我们信仰的人相同的异议），他凭自己的权柄责备他们无知的大胆，说道：\BibleQuote{但有人要说：死人将怎样复活呢？他们将带着什么样的身体回来呢？糊涂人哪！你所播的种子，若不先死去，决不会生发出来；并且你所播的，并不是那将要生出的形体，而是一颗赤裸的籽粒，譬如麦粒，或别的种粒；但天主随自己的心意给它一个形体。}在我看来，在这段话中，他堵住了那些人的口，这些人表现出对自然界适当比例的无知，并以自己的力量衡量天主的能力，认为只有人类智慧所能理解的范围才对天主是可能的，而超出此范围的也超出了天主的能力。因为那个曾问宗徒\Quote{死人怎样复活？}的人，显然暗示一旦身体的原子散开，它们就不可能再次聚合；而既然这是不可能的，并且除了由这种聚合所产生的身体之外，没有其他可能的身体形态，他便像聪明的辩论家那样，通过一种三段论，得出他想证明的结论：如果身体是原子的集合，而这些原子的第二次集合是不可能的，那么复活的人将得到什么样的身体呢？这个结论似乎隐藏在前提的巧妙设计中，宗徒称之为\Quote{糊涂}，因为它来自那些未能从受造界的其他部分领悟天主能力之高明的人。因为，他略去了天主手中更崇高的奇迹（借由这些奇迹本可轻易使听众陷入两难，例如他本可以问\Quote{天体如何或从何而来，例如太阳、月亮或星座中所见者；穹苍、空气、水、大地从何而来？}），反而通过与我们一同生长且为人熟知的事物，使反对者显得不顾虑。\Quote{难道连农事不也教导你们，}他问，\Quote{那计算天主超越性能力却以自身能力为限的人是糊涂的吗？}种子从何处获得从它们生出的身体？这生长之前发生了什么？难道不是先有死亡吗？至少，如果一个紧密结合的整体的解体就是死亡的话；因为确实不能认为种子会生发出嫩芽，除非它在土壤中溶解，变得如此海绵状和多孔，以至将其自身的品质与土壤的邻近水分混合，从而转变为根和芽；甚至不止于此，而是再次变化为茎，茎上有节，像许多扣环一样束住它，使其能够挺直地承载着满载谷粒的穗。那么，在谷粒在土壤中溶解之前，所有这些属于谷粒的东西在哪里？然而，这个结果却来自那谷粒；如果那谷粒不曾先存在，穗就不会产生。因此，正如穗的\Quote{身体}从种子中出现，是天主巧妙的权能从那单一之物中产生了一切，并且它既不完全与那种子相同，也不完全相异，同样（她坚持说），借由这些在种子上发生的奇迹，你现在可以解释复活的奥迹。天主的能力，以其全能之丰盈，不仅恢复你那曾解体的身体，还为其增添伟大而华丽的成分，借此，人被以一种更加宏伟的方式装备。\par

\Quote{它种在可朽坏中，}他说，\Quote{它复活在不可朽坏中；它种在软弱中，它复活在能力中；它种在羞辱中，它复活在光荣中；它种在属生理的身体，它复活在属神的身体。}麦粒在土壤中分解后，留下了其体积的微小和形状的特殊性质，然而它并没有离开和失去自身，而是仍然以自身为中心，生长成为麦穗，尽管它在许多方面（如大小、光辉、复杂性、形状）超越了自己。同样，人类在死亡中放下了从情欲倾向中获得的所有那些特殊环境；即羞辱、可朽坏、软弱和年岁的特征；然而人类并没有丧失自身。它仿佛变成了麦穗；即进入不可朽坏、光荣、尊贵、能力和绝对的完美；进入一种状态，在其中其生命不再以纯本性的特有方式进行，而是进入了属神的、无情感的存在。因为属生理的身体的特点是总是处于流动之中，总是从当前状态变化并变成别的东西；但这些过程——我们不仅在人类中观察到，也在植物和野兽中观察到——在将来的生命中都不会存留。此外，在我看来，宗徒的话完全符合我们自己对复活的理解。它们表明了我们自己在复活定义中所包含的同一件事，即我们曾说复活无非是\Quote{\Emph{我们本性在其原始形态中的重建}}。因为，我们从圣经中关于最初创造的记述中得知，首先大地生出了\BibleQuote{青草}（如叙述所说），然后从这植物结出了种子，当种子落在地上时，又生出与原来植物相同形状的植物；应当注意，宗徒宣称在复活中也发生同样的事；因此我们从宗徒那里学到一点，即我们的身体不仅将变成更高贵的东西，而且我们在其中所期待的不过是起初就有的东西。我们看到，起初，不是麦穗从麦粒生出，而是麦粒来自麦穗，然后麦穗围绕麦粒生长：所以这个比喻所显示的次序清楚地表明，通过复活为我们产生的那整个有福状态，只是回到我们原始的恩宠状态。事实上，我们曾经在某种意义上是一个丰满的麦穗；但罪的热焰使我们枯干，然后我们在死亡中分解，大地接受了我们；但在复活的春天，大地将重新生出我们身体这赤裸的麦粒，使其成为麦穗的形状，高大、匀称、挺拔，直达天顶，并且有叶子、有须、闪耀着不可朽坏的光辉，以及一切其他肖似天主的标记。因为\BibleQuote{这可朽坏的必须穿上不可朽坏}；而这不可朽坏、光荣、尊贵和能力正是那些与众不同的、公认的天主性标记，它们曾经属于按天主的肖像所造的人，并且我们期望在将来得到它们。第一个人亚当就是第一个麦穗；但是随着邪恶的到来，人性缩减为单纯的众多；并且，正如发生在麦穗上的麦粒一样，每一个人都被剥夺了那原始麦穗的美，并在土壤中腐朽；但在复活中，我们重生于我们原初的光辉；只是代替那个单一的原始麦穗，我们成为田野中无数麦穗的万万。与邪恶生活相对的美德生活是这样区分的：那些在活着时通过有德行的行为对自己进行耕种的人，立刻在完美的麦穗的一切品质中显现出来；而那些其赤裸的麦粒（即他们自然灵魂的力量）因邪恶的习惯而退化了，并且像被气候硬化了（如那些懂行的人所说的所谓\Quote{\OriginalTerm{角击}{hornstruck}}种子那样生长）的人，虽然他们将在复活中重生，但将从他们的审判者那里经历极大的严厉，因为他们没有力量成长为完整的麦穗，从而成为我们在地上堕落之前的样子。庄稼的监管者提供的补救方法是收集莠子和荆棘，这些莠子和荆棘与好种子一同生长，且所有曾经滋养其根的隐藏力量都进入了它们私生子的生命，以致它不仅不得不没有营养而存留，而且被这种不自然的生长所窒息，因而变得不结果实。当从它们内部的营养部分中，所有与之相反或伪造的东西都被拣出，并交付给吞噬一切不自然之物的火，然后消失，那么在这类人身上，他们的人性也将因这种耕作而繁荣并成熟为结果实，并在漫长的岁月之后，最终重新获得天主起初印在我们身上的那普遍形式。确实，那些在复活中一出生就发育出那些麦穗的丰满之美的人是有福的。然而，我们这样说，并不意味着暗示那些在今生度美德生活的人和度邪恶生活的人之间会有任何纯粹身体上的区别，好像我们应该认为一个人在其物质框架方面是不完美的，而另一个人在此方面会获得完美。囚徒和自由人，在此现世，就他们两人的身体结构而言是完全一样的；虽然就享受和痛苦而言，他们之间有巨大的鸿沟。依我看，我们应该以这种方式来估计善人与恶人在那个中间时期的区别。因为从死亡那种播种中复活的身体的完美，如宗徒告诉我们，在于不可朽坏、光荣、尊贵和能力；但任何这些卓越之处的减少，并不表示复活者相应身体上的残缺，而是从那些被认为属于善的事物中的每一种的退隐和隔绝。因此，既然这两种完全相反的意念（我指善与恶）中的一种必然以某种方式附着于我们，那么说一个人不包括在善之中，就是必要的证明他包括在恶之中。但是，在与恶的联系中，我们找不到尊贵、光荣、不可朽坏、能力；因此我们不得不打消一切疑虑：一个与这些事物毫无关系的人，必然与它们的对立面有关，即软弱、羞辱、可朽坏以及一切此类性质的事物，正如我们在讨论的前面部分提到的，当我们说有多少由恶生出的情欲，当它们渗透到整个自然的实质并与之成为一体时，灵魂多么难以摆脱它们。当这些情欲被火所完成的治愈过程净化并从其中彻底除去时，那么构成我们关于善之观念的一切事物都会取而代之；即不可朽坏、生命、尊贵、恩宠、光荣，以及我们猜想在天主内和按天主的肖像所造的人中将看到的一切其他事物。\par

\SourceRecord{Translated by William Moore and Henry Austin Wilson.；Nicene and Post-Nicene Fathers, Second Series；Vol. 5.；Edited by Philip Schaff and Henry Wace.；Buffalo, NY: Christian Literature Publishing Co.,；1893.；<http://www.newadvent.org/fathers/2915.htm>.}
